\section{Boundary Conditions}\label{SEC:BC}

\begingroup

\NewDocumentCommand{\smallSub}{m}{\ensuremath{\textsc{#1}}}
\NewDocumentCommand{\subBC}{}{\smallSub{bc}}
\NewDocumentCommand{\subL}{}{\smallSub{l}}
\NewDocumentCommand{\subR}{}{\smallSub{r}}
\NewDocumentCommand{\bcgroup}{m}{\textbf{#1}}
\NewDocumentCommand{\cvVecL}{}{\cvVec[\textsc{l}]}
\NewDocumentCommand{\cvVecR}{}{\cvVec[\textsc{r}]}
\NewDocumentCommand{\cvVecBC}{}{\cvVec[\textsc{bc}]}
\NewDocumentCommand{\dcnorm}{}{\ensuremath{\mathcal{B}}}

\subsection{How to Specify Boundary Conditions}\label{SEC:BC:Specify}

\subsubsection{Boundary Condition Values for
               \texorpdfstring{\inptext{Set\_BC\_Value}}%
                              {Set\_BC\_Value}}
\label{SEC:BC:Specify:SetBCValue}

%%%%%%%%%%%%%%%%%%%%%%%%%%%%%%%%%%%%%%%%%%%%%%%%%%%%%%%%%%%%%%%%%%%%%%%%%%%%%%%%

\devNotes{Need to add better error checking between input and CGNS grid file BC
names. Mike was seeing an error when it tries to sort the BC faces simply
because his input had \quoterm{Outlet} and the CGNS grid file had the name for
that BC as \quoterm{Outflow}. The error occurring while sorting the BC faces
provided ZERO indication that it had something to do with something as trivial
as a mismatch for the BC name between the input file and CGNS grid file.}

%%%%%%%%%%%%%%%%%%%%%%%%%%%%%%%%%%%%%%%%%%%%%%%%%%%%%%%%%%%%%%%%%%%%%%%%%%%%%%%%

The primary method for specifying which boundary condition to use is by
providing a value to the input
\BCinpBC[nmlBC::BCInput/SetBCValue][i]{Set\_BC\_Value},
where $i$ is an integer for the current boundary condition.
%
The supported values and their respective physical boundary conditions
are listed in \tref{TBL:BC:values}.
%
\begin{table}[h!]
  \begin{threeparttable}
% \caption{Types of boundary conditions available in GFR and their respective
%          input values to use with \BCinptext{Set\_BC\_Value}.}
  \caption{GFR boundary conditions.}
  \label{TBL:BC:values}
\begin{tabular}{ @{\quad} r @{\quad} p{0.7\textwidth} @{} }
  \toprule
  \inptext{Set\_BC\_Value} : & Type of Boundary Condition \\
 %\\[-1.25em] seems to be almost exactly a "Fortran backspace" equivalent
  \midrule\\[-0.75em]
  %-----------------------------------------------------------------------------
       & \bcgroup{Inflow boundary conditions}\\
   1 : & General inflow
           (\ref{SEC:BC:Inflow:General})\\
   2 : & Subsonic inflow
           (\ref{SEC:BC:Inflow:Subsonic})\\
   3 : & Supersonic inflow
           (\ref{SEC:BC:Inflow:Supersonic})\\
   5 : & SEM (unsteady turbulent) inflow%
           \tnote{\ensuremath{\ast}} \ %
           (\ref{SEC:BC:Inflow:SEM})\\[0.5em]
  %-----------------------------------------------------------------------------
       & \bcgroup{Outflow boundary conditions}\\
  11 : & General outflow
           (\ref{SEC:BC:Outflow:General})\\
  12 : & Subsonic outflow
           (\ref{SEC:BC:Outflow:Subsonic})\\
  13 : & Supersonic outflow
           (\ref{SEC:BC:Outflow:Supersonic})\\[0.5em]
  %-----------------------------------------------------------------------------
       & \bcgroup{Farfield boundary conditions}\\
  21 : & Characteristic \slash{} Riemann-invariant inflow\slash outflow%
           \tnote{\ensuremath{\ast}} \ %
           (\ref{SEC:BC:Farfield:Characteristic})\\
  22 : & General free-flow inflow\slash outflow
           (\ref{SEC:BC:Farfield:General})\\
  23 : & Freestream (fixed) farfield
           (\ref{SEC:BC:Farfield:Freestream})\\
  24 : & Fixed farfield
           (\ref{SEC:BC:Farfield:Fixed})\\[0.5em]
  %-----------------------------------------------------------------------------
       & \bcgroup{Wall boundary conditions}\\
  50 : & Inviscid\slash slip wall
           (\ref{SEC:BC:Wall:Slip})\\
  52 : & Adiabatic viscous\slash no-slip wall
           (\ref{SEC:BC:Wall:Adiabatic})\\
  53 : & Isothermal viscous\slash no-slip wall
           (\ref{SEC:BC:Wall:Isothermal})\\[0.5em]
  %-----------------------------------------------------------------------------
       & \bcgroup{Numerical boundary conditions}\\
  80 : & Symmetry boundary
           (\ref{SEC:BC:Numerical:Symmetry})\\
  81 : & Periodic (translated \andor rotated) coupled boundaries
          %(\ref{SEC:BC:Numerical:Periodic})\\
          %(\ref{SEC:BC:Numerical:Periodic})\\[0.50em] % w/o midrule
           (\ref{SEC:BC:Numerical:Periodic})\\[0.25em] % w/  midrule
  %-----------------------------------------------------------------------------
 %\midrule\\[-0.75em]
  \midrule\\[-1.0em]
  %-----------------------------------------------------------------------------
      %& \multicolumn{1}{@{}>{\footnotesize}p{0.75\textwidth}}{%
       & %\multicolumn{1}{@{}>{\small}p{0.75\textwidth}}{%
         \bcgroup{%
%        Other BC values used internally, just in case you encounter them.\newline
%        Note: \inptext{Set\_BC\_Value} should \emph{NOT} be set to any of
%        these values.
         %
         Other internal BC values that may be reported.\newline
         Note: Do \underline{NOT} set \inptext{Set\_BC\_Value} to any of these
         values.
         } %}
         \\
  %-----------------------------------------------------------------------------
  %-----------------------------------------------------------------------------
      0 : & Not a boundary condition\\
     99 : & Partition\slash communication boundary condition \\
  -9998 : & Unsupported boundary condition\\
  -9999 : & Unknown boundary condition\\
  %-----------------------------------------------------------------------------
  \bottomrule
\end{tabular}
  \begin{tablenotes}
    \item[\ensuremath{\ast}] Not compatible with multi-component gas mixtures.
  \end{tablenotes}
\end{threeparttable}
\end{table}

%%%%%%%%%%%%%%%%%%%%%%%%%%%%%%%%%%%%%%%%%%%%%%%%%%%%%%%%%%%%%%%%%%%%%%%%%%%%%%%%

As shown in \sref{SEC:Intro:FR}, the FR method needs to compute a common flux
(\ref{SEC:Intro:FR:Upwind}) at each interface flux point in order to create the
continuous flux function
(\ref{SEC:Intro:FR:ContinuousFlux}--\ref{SEC:Intro:FR:DerivativeFlux}) needed to
update the solution at the end of each time step
(\ref{SEC:Intro:FR:Time_Advancement}).
%
The inviscid common flux is computed using a Riemann solver (see
\inp[nmlNumerics::InvsFluxMethod]{Invs\_Flux\_Method}) which requires as input a
left and right state, \cvVecL{} and \cvVecR, respectively.
%
Furthermore, the viscous common flux requires the gradient of the left and right
states, $\nabla \cvVecL$ and $\nabla \cvVecR$, respectively.
%
See section \textit{II.B.2} in \linecite{Spiegel_2016} for details on the
discretization of the diffusion terms, and see inputs
\inp[nmlNumerics::ViscFluxMethod]{Visc\_Flux\_Method} and
\inp[nmlNumerics::ViscVariablesToAve]{Visc\_Variables\_To\_Ave}
on usage options in GFR.

%%%%%%%%%%%%%%%%%%%%%%%%%%%%%%%%%%%%%%%%%%%%%%%%%%%%%%%%%%%%%%%%%%%%%%%%%%%%%%%%

Due to the compact nature of FR\slash DG, each boundary condition really just
determines how the boundary state is evaluated before computing the common flux
function at the flux points of the respective boundary face.
%
Let \fupw[\textsc{i}]{} be the inviscid common flux function --- in this case
some kind of upwind or Riemann flux function --- defined as $\fupw[\textsc{i}] =
\fupw[\textsc{i}] \left( \cvVecL , \cvVecR \right)$.
%
Similarly, let \fupw[\textsc{v}]{} be the viscous common flux function defined
as $\fupw[\textsc{v}] = \fupw[\textsc{v}] \left( \cvVecL,\cvVecR ,
\nabla\cvVecL,\nabla\cvVecR \right)$.
%
For boundary faces, the left state is always from the interior cell and the
right state is the exterior state associated with the boundary condition.
%
Additionally, let \cvVecBC{} represent the boundary condition state, which is
specified in the input file for the boundary condition; \cvVecBC{} is related
--- but not necessary equivalent --- to \cvVecR.
%
In general, the right state is a function of both the left\slash interior state
and the boundary condition state, \thatis $\cvVecR = \cvVecR \left( \cvVecL ,
\cvVecBC \right)$.

%%%%%%%%%%%%%%%%%%%%%%%%%%%%%%%%%%%%%%%%%%%%%%%%%%%%%%%%%%%%%%%%%%%%%%%%%%%%%%%%

Sections \ref{SEC:BC:Inflow}--\ref{SEC:BC:Numerical} will provide details for
each boundary condition option currently available in GFR.
%
For each boundary condition type, this will cover what components of
\BCinpBC[nmlBC::BCInput]{} are required \andor optional to create the
%\hyperref[LST:nml_bc]{\bcInp{}} are required \andor optional to create the
needed information in \cvVecBC, and how that information and the interior state,
\cvVecL, are used to determine the right state, \cvVecR.
%
For most of the boundary conditions, the right gradient state is just copied
from the interior gradient state, \thatis
%
\begin{equation}
  \nabla\cvVecR=\nabla\cvVecL
  \label{EQ:BC:FlowThroughGradient}
\end{equation}
%
so this will be assumed for each boundary condition below unless otherwise
stated.

%%%%%%%%%%%%%%%%%%%%%%%%%%%%%%%%%%%%%%%%%%%%%%%%%%%%%%%%%%%%%%%%%%%%%%%%%%%%%%%%

Finally, for a given boundary condition, if both \BCinptext{Set\_BC\_Value} is
not set and the grid file does not provide a boundary condition type, the name
of the boundary condition --- taken from \BCinptext{Name} if provided, otherwise
taken from the name assigned to the boundary condition in the grid file ---
is parsed for certain keywords.
%
If a keyword match is found, a boundary condition will be assigned based on the
keyword.
%
Details on this are provided in \sref{SEC:BC:NameMatch}.

%%%%%%%%%%%%%%%%%%%%%%%%%%%%%%%%%%%%%%%%%%%%%%%%%%%%%%%%%%%%%%%%%%%%%%%%%%%%%%%%

\subsubsection{Fully Defining the Boundary Condition State}
\label{SEC:BC:Specify:FullState}

All flow-through boundary conditions (\thatis inflow and outflow) require a
minimum set of flow conditions to be defined in the input file in order to fully
define the boundary condition state.
%
There is one important detail to mention in how the boundary conditions are
processed before continuing on about what inputs are needed to fully define the
boundary condition state.
%
The first thing GFR does when processing the boundary conditions specified in
the input file is check if any individual boundary condition contains
all default flow conditions.
%
This check for default flow conditions checks all components of
\BCinpBC[nmlBC::BCInput]{}
\emph{except}:
%
\begin{InputExample}
  \item[\BCinptext{Name}]
  \item[\BCinptext{Set\_BC\_Value}]
  \item[\BCinptext{Periodic\%\ldots}]
\end{InputExample}
%
If an individual boundary condition is found to be set to the default flow
conditions, the flow state for that boundary condition will be set to the
freestream reference conditions specified through the
\hyperref[LST:nml_reference]{nml\_reference} namelist.
%
If just \emph{one} of \emph{any} of the other components (or sub-components) of
\BCinptext{} are \emph{not} set to the default value, the remainder of this
section on fully defining the boundary condition state applies to the
respective boundary condition.

\devNotes{If default flow conditions are found for an isothermal wall, it will
use the \texttt{ovar\_mod} variable \texttt{iso\_wall\_temp} to set the wall
temperature to 300 K. This probably needs to be changed to an error+abort. This
dev-note is also in \sref{SEC:BC:Wall:Isothermal}.}

%%%%%%%%%%%%%%%%%%%%%%%%%%%%%%%%%%%%%%%%%%%%%%%%%%%%%%%%%%%%%%%%%%%%%%%%%%%%%%%%

As discussed in \sref{SEC:nml_bc}, a variety of input variables are available
to make specifying the boundary flow conditions as easy as possible.
%
These input flow conditions include:
%
\begin{multicols}{2}
\begin{InputExample}
  \item[{\BCinpBC[nmlBC::BCInput/Pressure]{Pressure}}]
  \item[{\BCinpBC[nmlBC::BCInput/Temperature]{Temperature}}]
  \item[{\BCinpBC[nmlBC::BCInput/Density]{Density}}]
  \item[{\BCinpBC[nmlBC::BCInput/Mach]{Mach}}]
  \item[{\BCinpBC[nmlBC::BCInput/AoAAlpha]{AoA\_Alpha}}]
  \item[{\BCinpBC[nmlBC::BCInput/AoABeta]{AoA\_Beta}}]
  %
  \item[{\BCinpBC[nmlBC::BCInput/TotalPressure]{Total\_Pressure}}]
  \item[{\BCinpBC[nmlBC::BCInput/TotalTemperature]{Total\_Temperature}}]
  \item[{\BCinpBC[nmlBC::BCInput/TotalDensity]{Total\_Density}}]
  \item[{\BCinpBC[nmlBC::BCInput/VelocityX]{VelocityX}}]
  \item[{\BCinpBC[nmlBC::BCInput/VelocityY]{VelocityY}}]
  \item[{\BCinpBC[nmlBC::BCInput/VelocityZ]{VelocityZ}}]
\end{InputExample}
\end{multicols}

%%%%%%%%%%%%%%%%%%%%%%%%%%%%%%%%%%%%%%%%%%%%%%%%%%%%%%%%%%%%%%%%%%%%%%%%%%%%%%%%

As a general rule, four variables are sufficient to fully define the boundary
condition state.
%
If the velocity boundary condition has zero angle of attack,
using the default value of zero for both angles of attack reduces the minimum
number of variables to three.
%
After reading in the specified variables for a boundary condition, GFR will
iterate to define the full boundary state from the specified variables.
%
This iteration will loop five times and use what variables have been previously
set to define any remaining undefined state variables.
%
After completing the five iterations, if
any state variables remain undefined or are inconsistent with what was specified
in the input file, GFR will abort execution and report an error including the
values of all the input and computed state variables to help the user fix any
errors in the input file.

%%%%%%%%%%%%%%%%%%%%%%%%%%%%%%%%%%%%%%%%%%%%%%%%%%%%%%%%%%%%%%%%%%%%%%%%%%%%%%%%

One thing to note is that some boundary conditions may not actually need a fully
defined boundary state to define the right state.
%
An example of this would be a subsonic outflow boundary condition that only
requires the exterior static pressure to compute the right state, and most CFD
codes will only require a single input value for this quantity.
%
However, at the moment, GFR will require additional information to fully define
the boundary condition state, even though the boundary condition itself will
completely ignore all the other state variables and only use static pressure to
compute the right state.

%%%%%%%%%%%%%%%%%%%%%%%%%%%%%%%%%%%%%%%%%%%%%%%%%%%%%%%%%%%%%%%%%%%%%%%%%%%%%%%%

An example of a minimal set of variables to define a boundary condition
would be:
%
\begin{InputExample}
  \item[\textcolor{Green}{\checkmark{} Fully defined boundary condition}]
  \item[\BCinptext{Mach}]
  \item[\BCinptext{Temperature}]
  \item[\BCinptext{Total\_Pressure}]
\end{InputExample}
%
From this information, and using the default zero angle of attack, the velocity
vector and remaining static\slash total density\slash pressure\slash temperature
variables can be defined to fully define the boundary condition state.
%
The following example would not constitute a minimal set of variables to fully
define a boundary condition:
%
\begin{InputExample}
  \item[\textcolor{Red}{$\bm{\mathsf{X}}$ Incomplete boundary condition}]
  \item[\BCinptext{Density}]
  \item[\BCinptext{Temperature}]
  \item[\BCinptext{Pressure}]
\end{InputExample}
%
This boundary condition cannot be fully defined because no velocity condition
can be deduced from the given quantities.
%
However, this would be fully defined if one of these inputs were replaced with
their respective total\slash stagnation value, since the Mach number for the
boundary condition could be determined from having both the static and total
conditions.

%###############################################################################

\subsubsection{Situational Recommendations}
\label{SEC:BC:Specify:Recommendations}

This section includes some recommendations on which boundary conditions to use
in certain situations.

\begin{itemize}
  %
  \item
    We recommend characteristic boundary conditions for
    farfield boundaries that are generally aligned with the mean flow.
    %
    The next best option is the general free-flow boundary
    condition to allow both inflow and outflow through the boundary.
    %
    The farfield and fixed boundary conditions should best be avoided since
    they are essentially supersonic inflow boundary conditions, which could
    potentially cause problems if the interior flow propagates to the farfield
    boundary and needs to exit the computational domain.
    %
  \item
    The characteristic boundary condition works fairly well as a general
    outflow boundary condition, but is not recommended for a subsonic inflow
    since it does not enforce total conditions.
    %
    Maintaining total conditions is often necessary for internal or confined
    flow problems.
    %
  \item
    If you are not entirely sure if an inflow or outflow boundary will be either
    fully subsonic or supersonic, then use the general versions
    of those boundary conditions since they are set up to use the \quoterm{best}
    method of the respective boundary condition.
    %
    However, if you do know that it will be entirely subsonic or supersonic,
    it is best to use that respective version of the boundary condition.
    %
    This will slightly reduce the computational cost of determining the boundary
    condition states since it eliminates the need to evaluate the Mach number of
    the flow normal to the boundary face in order to determine which boundary
    condition to use.
    %
    The overall computational savings may
    be small and only noticeable over tens of thousands of time steps.
    %
  \item
    Symmetry boundary conditions should generally be avoided for LES since
    turbulent flows are never symmetric about a plane.
    %
    However, a slip wall boundary condition --- which is mathematically
    identical to a symmetry boundary conditions --- can be used for LES provided
    its usage \quoterm{makes physical sense}.
    %
    As an example, the boundary conditions used for the THX3 simulations (see
    \fref{FIG:Files:PW:THX3}) used a slip wall for the bottom boundary which
    included all surfaces on the same plane as the flat plate behind the square
    nozzle.
    %
    It might have been better to use a characteristic boundary condition to
    allow for downward spillage of the flow at the sides of the plate, but the
    slip walls seem to work sufficiently for this problem.%
    %
  \item
    Caution should be exercised when using periodic boundary conditions in the
    cross-stream directions for LES, since a narrow physical domain may limit
    the size of turbulent structures.
  %
\end{itemize}

%###############################################################################

\subsubsection{Setting Boundary Condition through Name Match}
\label{SEC:BC:NameMatch}

\NOTE The method described here for assigning a boundary condition type should
be avoided.
%
Instead, \BCinpBC[nmlBC::BCInput/SetBCValue]{Set\_BC\_Value} should be used to
specify the boundary condition type for all boundary conditions.

%%%%%%%%%%%%%%%%%%%%%%%%%%%%%%%%%%%%%%%%%%%%%%%%%%%%%%%%%%%%%%%%%%%%%%%%%%%%%%%%

For a given boundary condition, if both \inptext{Set\_BC\_Value} is not set and
a boundary condition type is not provided by the grid file, GFR will attempt to
assign a boundary condition type based on the boundary condition name.
%
This was the previous method used by GFR before \BCinptext{} was added to the
code.
%
It was retained for backwards compatibility mainly for use with Plot3D grid
files, but was included as a double redundancy for other grid files.

%%%%%%%%%%%%%%%%%%%%%%%%%%%%%%%%%%%%%%%%%%%%%%%%%%%%%%%%%%%%%%%%%%%%%%%%%%%%%%%%

Before continuing, it should be noted that it is \emph{highly recommended} that
all boundary condition types are specified in the input file using
\inptext{Set\_BC\_Value}.
%
If using a CGNS grid file, it is also good practice that each boundary condition
is given a GFR-compatible boundary condition type in the grid file as backup.
%
As stated at the beginning of this section, this method of using the boundary
condition name to assign a boundary condition type should be avoided.
%
It is only documented here to make the user aware that this behavior exists
within the code.

%%%%%%%%%%%%%%%%%%%%%%%%%%%%%%%%%%%%%%%%%%%%%%%%%%%%%%%%%%%%%%%%%%%%%%%%%%%%%%%%

The name for the boundary condition is taken from \BCinptext{Name} if provided;
otherwise, it is taken from the name assigned to the boundary condition in the
grid file.
%
The boundary condition name is parsed for keywords which will assign
certain boundary conditions depending on any keywords that might be found.
%
Before checking for a name match, the boundary condition name is processed
through the following steps to create a final character string for matching:
%
\begin{itemize}
  \item All alphabet characters are converted to upper case.
  \item All spaces, underscores, and hyphens are removed.
  \item The substring \quotett{SUPERSONIC} is reduced to \quotett{SUP}.
  \item The substrings \quotett{BOUNDARY}, \quotett{BC}, and \quotett{SONIC} are
    removed.
\end{itemize}
%
At this point, two special checks are attempted on the remaining string.
%
If any substring within the remaining string matches \quotett{SEM} or
\quotett{PERIODIC}, \inptext{Set\_BC\_Value} will be set to an SEM boundary
condition or a periodic boundary condition, respectively.
%
If neither of those substrings exist, then the remaining strings from
\tref{TBL:BC:NameMatch} will be used to provide a positive name match and
successfully assign a boundary condition type.
%
The first column of \tref{TBL:BC:NameMatch} lists the boundary condition value
assigned to \inptext{Set\_BC\_Value} if, after manipulating \inptext{Name}
using the above steps, the string that remains is equivalent to any of the
strings in the last column.
%
If no name match is made, \inptext{Set\_BC\_Value} is set to \texttt{-9999}
(Unknown boundary condition) and GFR will stop with an error message.

%%%%%%%%%%%%%%%%%%%%%%%%%%%%%%%%%%%%%%%%%%%%%%%%%%%%%%%%%%%%%%%%%%%%%%%%%%%%%%%%

\begin{table}[tb]
  \centering
% \caption{List of strings that will provide a positive match when trying to
%          assign the type of boundary condition based on a name match.}
  \caption{Boundary condition name strings.}
  \label{TBL:BC:NameMatch}
  \small{
  \begin{tabular}{ @{} r l @{\quad} c @{\qquad} l @{} }
    %%%%%%%%%%%%%%%%%%%%%%%%%%%%%%%%%%%%%%%%%%%%%%%%%%%%%%%%%%%%%%%%%%%%%%%%%%%%
    \toprule
    %%%%%%%%%%%%%%%%%%%%%%%%%%%%%%%%%%%%%%%%%%%%%%%%%%%%%%%%%%%%%%%%%%%%%%%%%%%%
   %\multirow[b]{2}{*}{Section \#} &
   %\multirow[b]{2}{*}{Boundary Condition Description} &
    \inptext{Set\_BC\_Value} & 
    Boundary Condition &
    Section \# &
    Matching Strings \\
   %(PW \texttt{"CAE Type"}) & Value & & \\ % Header bottom row
    %%%%%%%%%%%%%%%%%%%%%%%%%%%%%%%%%%%%%%%%%%%%%%%%%%%%%%%%%%%%%%%%%%%%%%%%%%%%
    \midrule
    %%%%%%%%%%%%%%%%%%%%%%%%%%%%%%%%%%%%%%%%%%%%%%%%%%%%%%%%%%%%%%%%%%%%%%%%%%%%
        1 & General Inflow
          & (\ref{SEC:BC:Inflow:General})
          & \quotett{INFLOW} \quotett{INLET} \\[0.5em]
    %%%%%%%%%%%%%%%%%%%%%%%%%%%%%%%%%%%%%%%%%%%%%%%%%%%%%%%%%%%%%%%%%%%%%%%%%%%%
        2 & Subsonic Inflow
          & (\ref{SEC:BC:Inflow:Subsonic})
          & \quotett{SUBINFLOW} \quotett{SUBIN} \\
      & & & \quotett{INFLOWSUB} \quotett{INSUB} \\[0.5em]
    %%%%%%%%%%%%%%%%%%%%%%%%%%%%%%%%%%%%%%%%%%%%%%%%%%%%%%%%%%%%%%%%%%%%%%%%%%%%
        3 & Supersonic Inflow
          & (\ref{SEC:BC:Inflow:Supersonic})
          & \quotett{SUPINFLOW} \quotett{SUPIN} \\
      & & & \quotett{INFLOWSUP} \quotett{INSUP} \\[0.5em]
    %%%%%%%%%%%%%%%%%%%%%%%%%%%%%%%%%%%%%%%%%%%%%%%%%%%%%%%%%%%%%%%%%%%%%%%%%%%%
        5 & SEM Turbulent Inflow
          & (\ref{SEC:BC:Inflow:SEM})
          & \quotett{SEM} \\[0.5em]
    %%%%%%%%%%%%%%%%%%%%%%%%%%%%%%%%%%%%%%%%%%%%%%%%%%%%%%%%%%%%%%%%%%%%%%%%%%%%
       11 & General Outflow
          & (\ref{SEC:BC:Outflow:General})
          & \quotett{OUTFLOW} \quotett{EXIT} \\[0.5em]
    %%%%%%%%%%%%%%%%%%%%%%%%%%%%%%%%%%%%%%%%%%%%%%%%%%%%%%%%%%%%%%%%%%%%%%%%%%%%
       12 & Subsonic Outflow
          & (\ref{SEC:BC:Outflow:Subsonic})
          & \quotett{SUBOUTFLOW} \quotett{SUBOUT} \\
      & & & \quotett{OUTFLOWSUB} \quotett{OUTSUB} \\[0.5em]
    %%%%%%%%%%%%%%%%%%%%%%%%%%%%%%%%%%%%%%%%%%%%%%%%%%%%%%%%%%%%%%%%%%%%%%%%%%%%
       13 & Supersonic Outflow
          & (\ref{SEC:BC:Outflow:Supersonic})
          & \quotett{SUPOUTFLOW} \quotett{SUPOUT} \\
      & & & \quotett{OUTFLOWSUP} \quotett{OUTSUP} \\
      & & & \quotett{EXTRAP} \quotett{EXTRAPOLATE} \\[0.5em]
    %%%%%%%%%%%%%%%%%%%%%%%%%%%%%%%%%%%%%%%%%%%%%%%%%%%%%%%%%%%%%%%%%%%%%%%%%%%%
       21 & Characteristic Inflow\slash Outflow
          & (\ref{SEC:BC:Farfield:Characteristic})
          & \quotett{CHARACTERISTIC} \quotett{CHAR} \quotett{CHARACTER}
          \\[0.5em]
    %%%%%%%%%%%%%%%%%%%%%%%%%%%%%%%%%%%%%%%%%%%%%%%%%%%%%%%%%%%%%%%%%%%%%%%%%%%%
       22 & General Inflow\slash Outflow
          & (\ref{SEC:BC:Farfield:General})
          & \quotett{GENERAL} \quotett{FREEFLOW} \quotett{FARFIELD} \\[0.5em]
    %%%%%%%%%%%%%%%%%%%%%%%%%%%%%%%%%%%%%%%%%%%%%%%%%%%%%%%%%%%%%%%%%%%%%%%%%%%%
       23 & Freestream (Fixed) Farfield
          & (\ref{SEC:BC:Farfield:General})
          & \quotett{FREESTREAM} \quotett{REFERENCE} \\[0.5em]
    %%%%%%%%%%%%%%%%%%%%%%%%%%%%%%%%%%%%%%%%%%%%%%%%%%%%%%%%%%%%%%%%%%%%%%%%%%%%
       24 & Fixed Farfield
          & (\ref{SEC:BC:Farfield:General})
          & \quotett{FIXED} \\[0.5em]
    %%%%%%%%%%%%%%%%%%%%%%%%%%%%%%%%%%%%%%%%%%%%%%%%%%%%%%%%%%%%%%%%%%%%%%%%%%%%
       50 & Slip Wall
          & (\ref{SEC:BC:Wall:Slip})
          & \quotett{SWALL} \quotett{SLIP} \quotett{SLIPWALL} \\
      & & & \quotett{WALLINVISCID} \\
      & & & \quotett{EWALL} \quotett{EULER} \quotett{EULERWALL} \\[0.5em]
    %%%%%%%%%%%%%%%%%%%%%%%%%%%%%%%%%%%%%%%%%%%%%%%%%%%%%%%%%%%%%%%%%%%%%%%%%%%%
       52 & Adiabatic Wall
          & (\ref{SEC:BC:Wall:Adiabatic})
          & \quotett{AWALL} \quotett{ADIABATIC} \quotett{WALL} \\
      & & & \quotett{WALLVISCOUS} \quotett{WALLVISCOUSHEATFLUX} \\[0.5em]
    %%%%%%%%%%%%%%%%%%%%%%%%%%%%%%%%%%%%%%%%%%%%%%%%%%%%%%%%%%%%%%%%%%%%%%%%%%%%
       53 & Isothermal Wall
          & (\ref{SEC:BC:Wall:Isothermal})
          & \quotett{IWALL} \quotett{ISOTHERMAL} \\
      & & & \quotett{WALLVISCOUSISOTHERMAL} \\[0.5em]
    %%%%%%%%%%%%%%%%%%%%%%%%%%%%%%%%%%%%%%%%%%%%%%%%%%%%%%%%%%%%%%%%%%%%%%%%%%%%
       80 & Symmetry Boundary
          & (\ref{SEC:BC:Numerical:Symmetry})
          & \quotett{SYMM} \quotett{SYMMETRY} \quotett{SYMMETRYPLANE} \\
    %%%%%%%%%%%%%%%%%%%%%%%%%%%%%%%%%%%%%%%%%%%%%%%%%%%%%%%%%%%%%%%%%%%%%%%%%%%%
       81 & Periodic Boundary
          & (\ref{SEC:BC:Numerical:Periodic})
          & \quotett{PERIODIC} \\
    \bottomrule
  \end{tabular}
  }
\end{table}

%%%%%%%%%%%%%%%%%%%%%%%%%%%%%%%%%%%%%%%%%%%%%%%%%%%%%%%%%%%%%%%%%%%%%%%%%%%%%%%%

\devNotes{Need to modify \inptext{bc\_string\_to\_integer} so that
\quotett{GENERAL} returns \quotett{bc\_generic\_freeflow}. Also need to add
matches for \inptext{SUBIN, INSUB, SUPIN, INSUP, SUBOUT, OUTSUB, SUPOUT,
OUTSUP}. Also need to add removal of hyphens.}

%%%%%%%%%%%%%%%%%%%%%%%%%%%%%%%%%%%%%%%%%%%%%%%%%%%%%%%%%%%%%%%%%%%%%%%%%%%%%%%%



%###############################################################################

\subsection{Inflow Boundaries}\label{SEC:BC:Inflow}

%###############################################################################

\subsubsection{General Inflow Boundary Condition}
\label{SEC:BC:Inflow:General}

This boundary condition combines subsonic and supersonic inflow boundary
conditions into a single option.
%
This boundary condition takes the velocity from the left\slash interior state,
$\vec{u}_\subL$, and evaluates the magnitude of that velocity normal to the
boundary face, $v_{n,\subL}$, as
%
\[
  v_{n,\subL} = \vec{u}_\subL\cdot\normal
\]
%
where \normal{} is the unit normal for the boundary face that points outside the
computational domain.
%
The interior static temperature $T_\subL$ and $v_{n,\subL}$ can then be used to
compute an interior Mach number normal to the inflow boundary, $M_{n,\subL}$, as
%
\[
  M_{n,\subL} = \frac{v_{n,\subL}}{\sqrt{\gamma \rgas T_\subL}}
\]
%
If $M_{n,\subL} \leq -1$ (negative because the normal points outside the
domain), the inflow is determined to be supersonic, and the right state is
evaluated the same way as the supersonic inflow of
\sref{SEC:BC:Inflow:Supersonic} by just fixing it to
the boundary condition state, \thatis \[ \cvVecR = \cvVecBC \]
%
Otherwise, the inflow is subsonic, and the right state is evaluated based on the
subsonic inflow boundary condition with $\inptext{Sub\_Inflow\_Method}=4$
(\ref{SEC:BC:Inflow:Subsonic:4}).

%%%%%%%%%%%%%%%%%%%%%%%%%%%%%%%%%%%%%%%%%%%%%%%%%%%%%%%%%%%%%%%%%%%%%%%%%%%%%%%%

It should be noted that this boundary condition will always operate as if
$\inptext{Sub\_Inflow\_Method}=4$ when the interior Mach number is found to be
subsonic.  The value of the input \inptext{Sub\_Inflow\_Method} has no control
over the subsonic inflow method used for this boundary condition and is
therefore ignored.

%%%%%%%%%%%%%%%%%%%%%%%%%%%%%%%%%%%%%%%%%%%%%%%%%%%%%%%%%%%%%%%%%%%%%%%%%%%%%%%%

\devNotes{Need to fix subsonic inflow angle-of-attack bug.}

%###############################################################################



%###############################################################################

\subsubsection{Subsonic Inflow Boundary Condition}
\label{SEC:BC:Inflow:Subsonic}

We first need to introduce notation for some quantities that will be utilized
frequently when describing the different subsonic inflow methods.
%
For the velocity vector $\vec{u}$, let \vmag{} represent the total magnitude of
the vector, \thatis
%
\[ \vmag = \norm{\vec{u}} \]
%
We also introduce the following notation to represent the vector of the
directional cosines for each respective component of a velocity vector
%
\begin{equation}
  \dircos = \frac{\vec{u}}{\vmag}
  \label{EQ:BC:Inflow:Subsonic:DirCosines}
\end{equation}
%
where the accent ${\hat{\phantom{a}}}$ is used here to denote this is a unit
vector by definition.
%
For most of the following subsonic inflow methods, the directional cosines of
the prescribed inflow velocity are used to fix the angle of attack for the right
state velocity, therefore $\dircos_\subR=\dircos_\subBC$.

%%%%%%%%%%%%%%%%%%%%%%%%%%%%%%%%%%%%%%%%%%%%%%%%%%%%%%%%%%%%%%%%%%%%%%%%%%%%%%%%

Using these definitions, the velocity in the direction of the
(outward pointing) surface normal can be written as
%
\begin{equation}
  v_n = \vec{u}\cdot\normal
      = \left(\vmag\dircos\right)\cdot\normal
      = \vmag\left(\dircos\cdot\normal\right)
      = \vmag\dcnorm
  \label{EQ:BC:Inflow:Subsonic:NormalVelocity}
\end{equation}
%
where
\vmag{} is the magnitude of the velocity vector defined above and
\dcnorm{} represents the fraction of \vmag{} in the direction of \normal.
%
For a flow that is perpendicular to the surface, $\dcnorm{} = -1$.
%
We also note that $v_n < 0$ on an inflow surface and
$\abs{v_n} < a$ on a subsonic inflow.


%%%%%%%%%%%%%%%%%%%%%%%%%%%%%%%%%%%%%%%%%%%%%%%%%%%%%%%%%%%%%%%%%%%%%%%%%%%%%%%%

The Riemann invariants and corresponding characteristic velocities are
given by
%
\begin{align}
      R^+       &= v_n + \frac{2a_\subL}{\gamma-1}
  &   \lambda_1 &= v_n + a > 0
  &&  \text{Outgoing}
  \\
      R^-       &= v_n - \frac{2a_\subL}{\gamma-1}
  &   \lambda_2 &= v_n - a < 0
  &&  \text{Incoming}
\end{align}

%%%%%%%%%%%%%%%%%%%%%%%%%%%%%%%%%%%%%%%%%%%%%%%%%%%%%%%%%%%%%%%%%%%%%%%%%%%%%%%%

It should be noted that for multi-component gas mixtures, the right state is
determined in a manner similar to method 4 (\ref{SEC:BC:Inflow:Subsonic:4})
unless method 2 (\ref{SEC:BC:Inflow:Subsonic:2}) is specified via
\inp[nmlBC::SubInflowMethod]{Sub\_Inflow\_Method}.

\devNotes{Need to fix subsonic inflow angle-of-attack bug for TPG subroutine.}

%%%%%%%%%%%%%%%%%%%%%%%%%%%%%%%%%%%%%%%%%%%%%%%%%%%%%%%%%%%%%%%%%%%%%%%%%%%%%%%%

\subsubsubsection{\texorpdfstring{\inptext{Sub\_Inflow\_Method} = 1}%
                                 {Sub\_Inflow\_Method = 1}}
\label{SEC:BC:Inflow:Subsonic:1}

This method performs a Newton iteration on the total temperature to get the
right state.
%
It is based on the method found on pages 8--9 of \linecite{Cooper_1990} from
the subsection titled \emph{Internal Flows} in Section 4.2.1.
%
This method uses the specified boundary condition values $p_{t,\subBC}$,
$T_{t,\subBC}$, and $\dircos_\subBC$ along with the outward propagating
Riemann invariant, $R^+$, to determine the right state.

%%%%%%%%%%%%%%%%%%%%%%%%%%%%%%%%%%%%%%%%%%%%%%%%%%%%%%%%%%%%%%%%%%%%%%%%%%%%%%%%

We start with enforcing that the total temperature for the right state is equal
to the prescribed total temperature
%
\[
  T_{t,\subR} = T_\subR \left(1 + \frac{\gamma-1}{2}M^2_\subR\right)
              = T_{t,\subBC}
\]
%
Using the relations $M=\vmag/a$, $a^2=\gamma\rgas T$, and
$c_p=\gamma\rgas/\left(\gamma-1\right)$, the total temperature can be rewritten
as
%
\begin{equation}
  T_{t,\subBC} = T_\subR + \frac{1}{2c_p}\vmag^2_\subR
  \label{EQ:BC:Inflow:Subsonic1:Ttotal}
\end{equation}
%
The next step is to set the outgoing Riemann invariant, $R^+$, using the
interior state.
%
We use the subscript $\subL$ to note that the invariant is computed from
the left (interior) state.
%
\begin{equation}
  R^+_\subL = \vec{u}_\subL\cdot\normal + \frac{2a_\subL}{\gamma-1}
  \label{EQ:BC:Inflow:Subsonic1:RplusL}
\end{equation}
%
We want to find a right state normal velocity and speed of sound that keeps the
outgoing invariant constant, or
%
\begin{equation}
  \vec{u}_\subR\cdot\normal + \frac{2a_\subR}{\gamma-1} = R^+_\subL
  \label{EQ:BC:Inflow:Subsonic1:Rplus}
\end{equation}
%
Using \eref{EQ:BC:Inflow:Subsonic:NormalVelocity} to represent the velocity
normal to the surface, we can rearrange \eref{EQ:BC:Inflow:Subsonic1:Rplus} as
%
\begin{equation}
  \vmag_\subR = \frac{1}{\dcnorm_\subBC}
                \left( R^+_\subL - \frac{2a_\subR}{\gamma-1} \right)
  \label{EQ:BC:Inflow:Subsonic1:vmagR}
\end{equation}
%
where the subscript \subBC{} has been added to $\dcnorm_\subBC$ to indicate
that we are using the direction cosines from the boundary condition velocity
vector to enforce the right state velocity to the prescribed angles of attack.
%
Finally, we can combine \eref{EQ:BC:Inflow:Subsonic1:Ttotal} and
\eref{EQ:BC:Inflow:Subsonic1:vmagR} resulting in the equation
%
\begin{equation*}
  T_{t,\subBC} = T_\subR + \frac{1}{2c_p}
    \left[ 
      \frac{1}{\dcnorm_\subBC}
      \left( R^+_\subL - \frac{2a_\subR}{\gamma-1} \right)
    \right]^2
\end{equation*}
%
With $a_\subR=\sqrt{\gamma\rgas T_\subR}$, this equation becomes solely
dependent on the right state temperature $(T_\subR)$ with all the other
variables being either constant or prescribed by the boundary condition state
%
\begin{equation}
  f\left(T_\subR\right) = T_\subR + \frac{1}{2c_p}
    \left[ 
      \frac{1}{\dcnorm_\subBC}
      \left( R^+_\subL - \frac{2a_\subR}{\gamma-1} \right)
    \right]^2 - T_{t,\subBC}
  = 0
  \label{EQ:BC:Inflow:Subsonic1:fTr}
\end{equation}
%
We can solve this numerically using a Newton iteration
%
\begin{align*}
  T^{n+1}_\subR =&
         T^n_\subR - \frac{f\left(T^n_\subR\right)}{f'\left(T^n_\subR\right)}
  \\
  a^{n+1}_\subR =& \sqrt{ \gamma \rgas T^{n+1}_\subR }
\end{align*}
%
where $n$ is the current iteration of the Newton solver.
%
All that remains is taking the derivative of $f\left(T_\subR\right)$, which
condenses nicely to
%
\begin{equation*}
  f'\left(T_\subR\right) = 1 - \frac{\vmag_\subR}{\dcnorm_\subBC a_\subR}
\end{equation*}
%
After the Newton iteration converges to a solution for $T_\subR$ and
$a_\subR$,
%the speed of sound is updated using $a_\subR=\sqrt{\gamma\rgas T_\subR}$,
we can then continue on to compute the final right state variables needed
for the boundary condition.
%
The next step is to update $\vmag_\subR$ using
\eref{EQ:BC:Inflow:Subsonic1:vmagR} from which a Mach number,
$M_\subR=\vmag_\subR / a_\subR$, can be evaluated.
%
Finally, the primitive variables for the right state can be computed as
%
\begin{alignat}{1}
  \vec{u}_\subR & = \vmag_\subR \dircos_\subBC
  \label{EQ:BC:Inflow:Subsonic1:velR}
  \\
        p_\subR & = p_{t,\subBC}
                    \left( 1 + \frac{\gamma-1}{2} M^2_\subR
                    \right)^{-\gamma / \left( \gamma - 1 \right)}
  \label{EQ:BC:Inflow:Subsonic1:presR}
  \\
     \rho_\subR & = \frac{p_\subR}{\rgas T_\subR}
  \label{EQ:BC:Inflow:Subsonic1:densR}
\end{alignat}

%%%%%%%%%%%%%%%%%%%%%%%%%%%%%%%%%%%%%%%%%%%%%%%%%%%%%%%%%%%%%%%%%%%%%%%%%%%%%%%%

\subsubsubsection{\texorpdfstring{\inptext{Sub\_Inflow\_Method} = 2}%
                                 {Sub\_Inflow\_Method = 2}}
\label{SEC:BC:Inflow:Subsonic:2}

This method is exactly the opposite of the subsonic outflow method, using the
density and velocity specified for the boundary condition and taking the static
pressure from the interior
%
\begin{alignat*}{1}
  \rho_\subR & = \rho_\subBC \\
  \vec{u}_\subR & = \vec{u}_\subBC \\
  p_\subR & = p_\subL
\end{alignat*}
%
This method does not enforce the total temperature or pressure at the inflow, so
caution should be taken with its use.
%
This method is probably only suitable for a subsonic inflow that is far upstream
from the body of interest. Even then, it is probably best to use one of the
other subsonic inflow methods or the characteristic boundary condition
(\ref{SEC:BC:Farfield:Characteristic}) instead.
%

%%%%%%%%%%%%%%%%%%%%%%%%%%%%%%%%%%%%%%%%%%%%%%%%%%%%%%%%%%%%%%%%%%%%%%%%%%%%%%%%

\subsubsubsection{\texorpdfstring{\inptext{Sub\_Inflow\_Method} = 3}%
                                 {Sub\_Inflow\_Method = 3}}
\label{SEC:BC:Inflow:Subsonic:3}

This method uses the specified boundary condition values $p_{t,\subBC}$,
$\dircos_\subBC$, and $H_{t,\subBC}$ (total enthalpy) along with the outward
propagating Riemann invariant, $R^+$, to determine the right state.
%
It is very similar to
$\inptext{Sub\_Inflow\_Method}=1$ (\ref{SEC:BC:Inflow:Subsonic:1}),
but a quadratic equation is solved for $a_\subR$
instead of using a Newton iteration for $T_\subR$.
%
It is also equivalent to the method found in Section 2.7 of
\linecite{Carlson_2011} (pages 7--9), though somewhat different notation is
used there and the flow is assumed to be normal to the boundary.

We start by rewriting \eref{EQ:BC:Inflow:Subsonic1:fTr} in terms of the
speed of sound, with $R^+$ computed from the left (interior) state using
\eref{EQ:BC:Inflow:Subsonic1:RplusL},
\begin{equation}
  \frac{a_\subR^2}{\gamma \rgas} + \frac{1}{2c_p}
    \left[ 
      \frac{1}{\dcnorm_\subBC}
      \left( R^+_\subL - \frac{2a_\subR}{\gamma-1} \right)
    \right]^2 - T_{t,\subBC}
  = 0
\end{equation}
%
and then rearranging the terms into a standard quadratic form
with $H_{t,\subBC} = c_p T_{t,\subBC}$.
\begin{equation}
    \left[ \dcnorm^2_\subBC + \frac{2}{\gamma-1} \right] a^2_\subR
  - \left[ 2 R^+_\subL \right] a_\subR
  + \frac{\gamma-1}{2}
    \left[ \left( R^+_\subL \right)^2 - 2 \dcnorm^2_\subBC H_{t,\subBC} \right]
  = 0
% \label{EQ:BC:Inflow:Subsonic3:Quadratic}
\end{equation}
%
Using the quadratic formula to solve this equation results in two solutions for
$a_\subR$.
%
If for some reason both solutions are negative, the speed of sound from the
interior state is used for the right state, \thatis $a_\subR = a_\subL$.
%
Otherwise $a_\subR$ is taken to be the maximum value between the two solutions.

%%%%%%%%%%%%%%%%%%%%%%%%%%%%%%%%%%%%%%%%%%%%%%%%%%%%%%%%%%%%%%%%%%%%%%%%%%%%%%%%

The temperature can be computed from $T_\subR = a_\subR^2 / \gamma \rgas$,
and the remaining primitive variables can be updated in the same way as in
\inptext{Sub\_Inflow\_Method}=1 (\sref{SEC:BC:Inflow:Subsonic:1}).
%
$\vmag_\subR$ is updated using \eref{EQ:BC:Inflow:Subsonic1:vmagR}
from which a Mach number, $M_\subR=\vmag_\subR / a_\subR$, can be evaluated.
%
Then $\vec{u}_\subR$, $p_\subR$, and $\rho_\subR$ are computed from
\eref{%
EQ:BC:Inflow:Subsonic1:velR,%
EQ:BC:Inflow:Subsonic1:presR,%
EQ:BC:Inflow:Subsonic1:densR%
}.

%%%%%%%%%%%%%%%%%%%%%%%%%%%%%%%%%%%%%%%%%%%%%%%%%%%%%%%%%%%%%%%%%%%%%%%%%%%%%%%%

\subsubsubsection{\texorpdfstring{\inptext{Sub\_Inflow\_Method} = 4}%
                                 {Sub\_Inflow\_Method = 4}}
\label{SEC:BC:Inflow:Subsonic:4}

This is the default method used for the subsonic inflow boundary condition,
because it is the simplest of the four available methods that also enforces the
total conditions specified for the boundary condition.
%
This method uses the specified boundary condition values $p_{t,\subBC}$,
$T_{t,\subBC}$, and $\dircos_\subBC$ along with the interior Mach number at the
boundary to determine the right state.
%
Since both $p_{t,\subBC}$ and $T_{t,\subBC}$ are specified, this implies that
$\rho_{t,\subBC}$ is also specified through the ideal gas law.

%%%%%%%%%%%%%%%%%%%%%%%%%%%%%%%%%%%%%%%%%%%%%%%%%%%%%%%%%%%%%%%%%%%%%%%%%%%%%%%%

It starts by enforcing that the right state Mach number is set to the interior
Mach number
\[
  M_\subR = M_\subL = \frac{\vmag_\subL}{a_\subL}
\]
%
which then allows us to compute the static temperature for the
right state based on the specified total temperature using
%
\[
  T_\subR = T_{t,\subBC} \left(1 + \frac{\gamma-1}{2}M^2_\subR\right)^{-1}
\]
%
From here, the magnitude of the right state velocity can be computed as
%
\[
  \vmag_\subR = M_\subR a_\subR = M_\subR \sqrt{\gamma \rgas T_\subR}
\]
%
which gives us everything that is needed to finally compute the primitive
variables for the right state
%
\begin{alignat}{1}
     \rho_\subR & = \left( \frac{p_{t,\subBC}}{\rgas T_{t,\subBC}} \right)
                    \left( 1 + \frac{\gamma-1}{2} M^2_\subR
                    \right)^{-1 / \left( \gamma - 1 \right)} \\
  \vec{u}_\subR & = \vmag_\subR \dircos_\subBC \\
        p_\subR & = \rho_\subR \rgas T_\subR
  \label{EQ:BC:Inflow:Subsonic4:PV2}
\end{alignat}

%%%%%%%%%%%%%%%%%%%%%%%%%%%%%%%%%%%%%%%%%%%%%%%%%%%%%%%%%%%%%%%%%%%%%%%%%%%%%%%%

%If looking at this from a practical stance separate of the GFR code, most users
%would want to specify a total pressure and total temperature for a subsonic
%inflow boundary condition.
%%
%In this case, you could think of this as first computing $p_\subR$ from
%\eref{EQ:BC:Inflow:Subsonic3:PV} using the boundary condition total pressure,
%$p_{t,\subBC}$, and then using the ideal gas law to compute $\rho_\subR$ from
%$p_\subR$ and $T_\subR$.
%%
%Either way, the right state variables in \eref{EQ:BC:Inflow:Subsonic4:PV}
%should come out the same assuming $\rho_{t,\subBC}$, $p_{t,\subBC}$, and
%$T_{t,\subBC}$ all conform to the ideal gas law.
%%
%Within GFR, this does not really matter since the requirement from
%\sref{SEC:BC:Specify:FullState} of having to fully define each flow-through
%boundary condition means these three total conditions are all known and
%consistent.

%%%%%%%%%%%%%%%%%%%%%%%%%%%%%%%%%%%%%%%%%%%%%%%%%%%%%%%%%%%%%%%%%%%%%%%%%%%%%%%%

%###############################################################################



%###############################################################################

\subsubsection{Supersonic Inflow Boundary Condition}
\label{SEC:BC:Inflow:Supersonic}

This is the simplest of the inflow boundary conditions.
%
Since the incoming flow is supersonic, no information from the interior of the
computational domain can propagate upstream to the inflow boundary condition.
%
Therefore, the right state is simply fixed to the boundary condition state,
\thatis
%
\begin{equation*}
  \cvVecR = \cvVecBC
\end{equation*}

%###############################################################################



%###############################################################################

\subsubsection{Synthetic Eddy Method (SEM) Boundary Condition}%
\label{SEC:BC:Inflow:SEM}

The Synthetic Eddy Method (SEM) provides an unsteady turbulent inflow boundary
condition that is crucial for accurate LES without requiring large run-up
sections of computational grid for fully developing a turbulent boundary layer
prior to the computational geometry of interest.
%
The user provides an inflow profile with the desired time-averaged flow
properties and turbulence statistics.
%
At each time step, the SEM will then generate an unsteady inflow profile with
turbulent fluctuations that, when time averaged, will reproduce the requested
distribution.
%
\Sref{SEC:SEM} provides extensive details on the synthetic eddy method:
%
\sref{SEC:SEM:theory} provides the theory behind the method as well as an
overview of how the method works;
%
\sref{SEC:SEM:usage} provides details on how to use the SEM in GFR including the
file format for the SEM profile, the relevant input variables with some
recommendations, and other aspects unique to the SEM implementation within GFR;
%
and \sref{SEC:SEM:FAQ} provides a collection of frequently asked questions.

%###############################################################################



%###############################################################################
%###############################################################################
%###############################################################################
%###############################################################################
%###############################################################################



%###############################################################################

\subsection{Outflow Boundaries}\label{SEC:BC:Outflow}

%###############################################################################

\subsubsection{General Outflow Boundary Condition}
\label{SEC:BC:Outflow:General}

This boundary condition combines subsonic and supersonic outflow boundary
conditions into a single option.
%
This boundary condition takes the velocity from the left\slash interior state,
$\vec{u}_\subL$, and evaluates the magnitude of that velocity normal to the
boundary face, $v_{n,\subL}$, as
%
\[
  v_{n,\subL} = \vec{u}_\subL\cdot\normal
\]
%
where \normal{} is the unit normal for the boundary face that points outside the
computational domain.
%
The interior static temperature $T_\subL$ and $v_{n,\subL}$ can then be used to
compute an interior Mach number normal to the outflow boundary, $M_{n,\subL}$, as
%
\[
  M_{n,\subL} = \frac{v_{n,\subL}}{\sqrt{\gamma \rgas T_\subL}}
\]
%
If $M_{n,\subL} \geq 1$, the outflow is supersonic, and the
right state is evaluated the same way as \sref{SEC:BC:Outflow:Supersonic} by
fixing it to the interior state, \thatis \[ \cvVecR = \cvVecL \]
%
Otherwise, the outflow is subsonic, and the right state is evaluated based on
the subsonic outflow boundary condition (\ref{SEC:BC:Outflow:Subsonic}), except
there is no safe guard for the outflow being supersonic since that situation is
already dealt with for this boundary condition; the safe guard on the flow
becoming reversed does still exist, but again, only provides minimal protection
against very minor violations of the assumed outflow.

%###############################################################################



%###############################################################################

\subsubsection{Subsonic Outflow Boundary Condition}%
\label{SEC:BC:Outflow:Subsonic}

This boundary condition assumes the velocity normal to the boundary face is
subsonic and leaving the computational domain.
%
The right state is computed as
%
\begin{alignat*}{1}
  \vec{u}_\subR & = \vec{u}_\subL \\
  T_\subR & = T_\subL \\
  p_\subR & = p_\subBC
\end{alignat*}
%
This boundary condition does have some built-in safe guards in case the
velocity at the boundary face becomes supersonic or reversed (\thatis inflow
instead of outflow).
%
However, these safe guards only provide minimal protection against
minor violations of the assumed subsonic outflow.
%
\begin{itemize}[itemsep=5pt,topsep=5pt,parsep=0pt]
  \item If you suspect your outflow boundary will likely contain large sections
    of both subsonic and supersonic outflow, you should consider using either
    the general outflow (\ref{SEC:BC:Outflow:General}) or characteristic
    (\ref{SEC:BC:Farfield:Characteristic}) boundary conditions.
  \item If you suspect there is a chance parts of your outflow boundary could
    reverse direction, at a minimum you should switch to using a characteristic
    (\ref{SEC:BC:Farfield:Characteristic}) boundary condition instead of a
    subsonic outflow boundary condition. However, the better option is to modify
    the grid so that the outflow boundary is in a region that maintains a flow
    directed out of the domain and is not susceptible to reversing that
    direction.
\end{itemize}

%###############################################################################



%###############################################################################

\subsubsection{Supersonic Outflow Boundary Condition}
\label{SEC:BC:Outflow:Supersonic}

This is the reverse of the supersonic inflow boundary condition.
%
Since the outgoing flow is supersonic, no information from the boundary
condition can propagate upstream into the computational domain.
%
Therefore, the right state is simply fixed to the interior\slash left state,
\thatis
%
\begin{equation*}
  \cvVecR = \cvVecL
\end{equation*}

%%%%%%%%%%%%%%%%%%%%%%%%%%%%%%%%%%%%%%%%%%%%%%%%%%%%%%%%%%%%%%%%%%%%%%%%%%%%%%%%

If the outflow is adjacent to a no-slip wall, this boundary condition could
result in stability problems within the subsonic region of the boundary layer.
%
If stability issues occur, it is recommended to switch to either a
characteristic (\ref{SEC:BC:Farfield:Characteristic}) or
general outflow (\ref{SEC:BC:Outflow:General}) boundary condition, since both of
these boundary conditions test the outgoing Mach number to determine whether to
use a subsonic or supersonic boundary condition.

%###############################################################################



%###############################################################################
%###############################################################################
%###############################################################################
%###############################################################################
%###############################################################################



%###############################################################################

\subsection{Farfield Boundaries}\label{SEC:BC:Farfield}

%###############################################################################

\subsubsection{Characteristic \slash{} Riemann-Invariant Boundary Condition}
\label{SEC:BC:Farfield:Characteristic}

The characteristic boundary condition is an inflow\slash outflow boundary
condition that works for both subsonic and supersonic flows.
%
This boundary condition uses the Riemann invariants to compute the right state
variables based on what flow regime the boundary condition state is in, which is
determined by the magnitude and direction normal to the boundary face of the
interior velocity.
%
Here are a few comments\slash warnings about using this boundary condition
before continuing on with describing how it is formulated:
%
\begin{itemize}[itemsep=5pt,topsep=5pt,parsep=0pt]
  \item When used as a subsonic inflow boundary condition, there is no
    guarantee that the total conditions at the inflow will be enforced.
    This should not be a problem if the inflow is significantly far
    upstream from the computational body of interest (\forexample external
    flows).
  \item This works well as a farfield boundary condition and has not shown any
    problems when the boundary condition velocity is aligned tangent to the
    boundary face.
  \item This works well for an outflow boundary condition for both internal and
    external flows, whether they are subsonic, supersonic, or mixed.
  \item If used for an outflow boundary condition, the prescribed velocity for
    the boundary condition is important for subsonic outflows, whereas the
    boundary condition velocity is ignored for the subsonic outflow boundary
    condition in \sref{SEC:BC:Outflow:Subsonic}. It should be noted that,
    similar to when used for subsonic inflows, this boundary condition applied
    to subsonic outflows does not guarantee that any specified boundary
    condition state will be held constant.
\end{itemize}

%%%%%%%%%%%%%%%%%%%%%%%%%%%%%%%%%%%%%%%%%%%%%%%%%%%%%%%%%%%%%%%%%%%%%%%%%%%%%%%%

Recalling that \normal{} is the outward-facing unit normal for the boundary
face, let $v_{n,\subL}=\vec{u}_\subL\cdot\normal$ be the magnitude of the
interior velocity normal to the boundary face, and
$v_{n,\subBC}=\vec{u}_\subBC\cdot\normal$ be the same but for the prescribed
boundary condition velocity.
%
Similarly, let $a_\subL$ and $a_\subBC$ respectively be the speed of sound for
the interior and boundary condition states.
%
In order to compute the final right state, we need to compute four quantities:
$R^+$, the outgoing Riemann invariant; $R^-$, the incoming Riemann invariant;
$\vec{v}_{t,\subR}$, the tangential velocity vector at the boundary
face; and $s_\subR$, the entropy for the right state.

%%%%%%%%%%%%%%%%%%%%%%%%%%%%%%%%%%%%%%%%%%%%%%%%%%%%%%%%%%%%%%%%%%%%%%%%%%%%%%%%

The outgoing Riemann invariant, $R^+$, is evaluated as
%
\[
  R^+ =
  \begin{cases}
    v_{n,\subBC} + \dfrac{2 a_\subBC}{\gamma-1} &
      \text{\;if } \dfrac{v_{n,\subBC}}{a_\subBC} \leq -1
    \quad \text{(supersonic inflow)} \\[2.5ex]
    v_{n,\subL} + \dfrac{2 a_\subL}{\gamma-1} &
      \text{\;otherwise}
  \end{cases}
\]
The incoming Riemann invariant, $R^-$, is evaluated as
%
\[
  R^- =
  \begin{cases}
    v_{n,\subL} - \dfrac{2 a_\subL}{\gamma-1} &
      \text{\;if } \dfrac{v_{n,\subL}}{a_\subL} \geq 1
    \quad \text{(supersonic outflow)} \\[2.5ex]
    v_{n,\subBC} - \dfrac{2 a_\subBC}{\gamma-1} &
      \text{\;otherwise}
  \end{cases}
\]
%
The tangential velocity vector, $\vec{v}_{t,\subR}$, is evaluated as
%
\[
  \vec{v}_{t,\subR} =
  \begin{cases}
    \vec{u}_\subBC - \left(\vec{u}_\subBC\cdot\normal\right)\normal &
      \text{\;if } v_{n,\subL} \leq 0 \quad \text{(inflow)} \\[1.5ex]
    \vec{u}_\subL - \left(\vec{u}_\subL\cdot\normal\right)\normal &
      \text{\;if } v_{n,\subL} > 0 \quad \text{(outflow)}
  \end{cases}
\]
%
Finally, the entropy, $s_\subR$, is evaluated as
\[
  s_\subR =
  \begin{cases}
    \dfrac{p_\subBC}{\rho^{\gamma}_{\subBC}} &
      \text{\;if } v_{n,\subL} \leq 0 \quad \text{(inflow)} \\[2.5ex]
    \dfrac{p_\subL}{\rho^{\gamma}_{\subL}} &
      \text{\;if } v_{n,\subL} > 0 \quad \text{(outflow)}
   %  \\[2.0ex]
   %p_\subBC / \rho^{\gamma}_{\subBC} &
   %  \text{\;if } v_{n,\subL} \leq 0 \quad \text{(inflow)} \\[1.0ex]
   %p_\subL / \left(\rho^{\gamma}_{\subL}\right) &
   %  \text{\;if } v_{n,\subL} > 0 \quad \text{(outflow)}
  \end{cases}
\]

%%%%%%%%%%%%%%%%%%%%%%%%%%%%%%%%%%%%%%%%%%%%%%%%%%%%%%%%%%%%%%%%%%%%%%%%%%%%%%%%

These four quantities are used to compute the right state in the following way.
%
First, the characteristic normal velocity is computed by averaging the Riemann
invariants
%
\[
  v_{n,\subR} = \frac{1}{2}\left(R^+ + R^-\right)
\]
%
From there, the right state velocity can be determined by simply adding the
normal and tangential components.
%
\[
  \vec{u}_\subR = v_{n,\subR}\normal + \vec{v}_{t,\subR}
\]
%
Next, the speed of sound is computed.
%
\[
  a_\subR = \frac{\gamma-1}{4} \left( R^+ - R^- \right)
\]
%
Finally, the right state density and pressure can be computed as
%
\begin{alignat*}{1}
  \rho_\subR &
      = \left(\frac{a^2_\subR}{\gamma s_\subR}\right)^{\tfrac{1}{\gamma-1}}
      \\[1.0ex]
  p_\subR &
      = \frac{a^2_\subR \rho_\subR}{\gamma}
\end{alignat*}

%%%%%%%%%%%%%%%%%%%%%%%%%%%%%%%%%%%%%%%%%%%%%%%%%%%%%%%%%%%%%%%%%%%%%%%%%%%%%%%%

\devNotes{Currently, the characteristic speed of sound is calculated as $a_c =
(R^+ - R^-)(\gamma-1)/4$. From the first paragraph on page 379 of Hirsch vol.
2, this does not enforce total temperature, and that the equation $a_c^2 =
(H_\subBC - \vec{u}^2_\subR)(\gamma-1)/2$ should be used instead. Also the way
pressure is computed should be changed from $p=s\rho^{\gamma}$ to $p=a_c^2
\rho/\gamma$. These are equivalent, but the latter is just nicer without the
$\gamma$ in the exponent.
%
When switching to this notation, need to clarify what the subscript $c$
means in $a_c$ and how it relates to $a_\subR$.
}

%###############################################################################



%###############################################################################

\subsubsection{General Free-flow Boundary Condition}
\label{SEC:BC:Farfield:General}

This boundary condition serves as a generalized inflow\slash outflow for both
subsonic and supersonic flow regimes, essentially combining the general inflow
(\ref{SEC:BC:Inflow:General}) and general outflow (\ref{SEC:BC:Outflow:General})
boundary conditions into a single option.
%
This boundary condition takes the velocity from the left\slash interior state,
$\vec{u}_\subL$, and evaluates the magnitude of that velocity normal to the
boundary face, $v_{n,\subL}$, as
%
\[
  v_{n,\subL} = \vec{u}_\subL\cdot\normal
\]
%
where \normal{} is the unit normal for the boundary face that points outside the
computational domain.
%
The interior static temperature $T_\subL$ and $v_{n,\subL}$ can then be used to
compute an interior Mach number normal to the inflow boundary, $M_{n,\subL}$, as
%
\[
  M_{n,\subL} = \frac{v_{n,\subL}}{\sqrt{\gamma \rgas T_\subL}}
\]
%
From here, there are four options based on the magnitude and sign of
$M_{n,\subL}$:
%
\begin{itemize}
  %
  \item $M_{n,\subL}\geq +1$\quad (\textit{Supersonic Outflow})

    The interior flow is supersonic and exiting the domain.
    %
    The right state is evaluated the same way as the supersonic outflow boundary
    condition (\ref{SEC:BC:Outflow:Supersonic}) by just fixing it to the
    interior state, \thatis \[ \cvVecR = \cvVecL \]
  %
  \item $M_{n,\subL}\leq -1$\quad (\textit{Supersonic Inflow})

    The interior flow is supersonic because the magnitude of the Mach number is
    greater than or equal to 1. It is entering the domain since a negative value
    indicates the flow is aligned in the opposite direction as the outward
    facing normal.
    %
    The right state is evaluated the same way as the supersonic inflow boundary
    condition (\ref{SEC:BC:Inflow:Supersonic}) by just fixing it to the boundary
    condition state, \thatis \[ \cvVecR = \cvVecBC \]
  %
  \item $0 \leq M_{n,\subL} < 1$\quad (\textit{Subsonic Outflow})

    The interior flow is subsonic and exiting the domain.
    %
    The right state is evaluated based on the subsonic outflow boundary
    condition (\ref{SEC:BC:Outflow:Subsonic}), except there are no safe guards
    in place for supersonic or reversed flows since the other three conditional
    statements of this bulleted list explicitly deal with those situations for
    this boundary condition.
  %
  \item $-1 < M_{n,\subL} < 0$\quad (\textit{Subsonic Inflow})

    The interior flow is subsonic because the magnitude of the Mach number is
    less than 1; it is entering the domain since a negative value indicates the
    flow is aligned in the opposite direction as the outward facing normal.
    %
    The right state is evaluated based on the subsonic inflow boundary condition
    with $\inptext{Sub\_Inflow\_Method}=4$ (\ref{SEC:BC:Inflow:Subsonic:4}).
    %
    It should be noted that this boundary condition will always operate as if
    $\inptext{Sub\_Inflow\_Method}=4$ when the interior Mach number indicates a
    subsonic inflow condition; the value of the input
    \inptext{Sub\_Inflow\_Method} has no control over the subsonic inflow
    method used for this boundary condition and is therefore ignored.
  %
\end{itemize}

%%%%%%%%%%%%%%%%%%%%%%%%%%%%%%%%%%%%%%%%%%%%%%%%%%%%%%%%%%%%%%%%%%%%%%%%%%%%%%%%

\devNotes{Need to fix subsonic inflow angle-of-attack bug.}

%###############################################################################



%###############################################################################

\subsubsection{Freestream Farfield Boundary Condition}%
\label{SEC:BC:Farfield:Freestream}

This is a supersonic inflow boundary condition, with the only difference between
this and the actual supersonic inflow boundary condition in
\sref{SEC:BC:Inflow:Supersonic} is this uses $Q_{\infty}$ from
the reference conditions in \sref{SEC:Reference_Conditions} as $Q_\subBC$.

\devNotes{Freestream boundaries are usually more flexible than this.
Consider making it more general to handle sub/super-sonic flow,
like the General Free-flow boundary condition, but using values from the
reference conditions.}

%###############################################################################



%###############################################################################

\subsubsection{Fixed Farfield Boundary Condition}\label{SEC:BC:Farfield:Fixed}

This is exactly the same thing as using a supersonic inflow boundary condition.
%
See \sref{SEC:BC:Inflow:Supersonic}.

%###############################################################################



%###############################################################################
%###############################################################################
%###############################################################################
%###############################################################################
%###############################################################################



%###############################################################################

\subsection{Wall Boundaries}\label{SEC:BC:Wall}

%###############################################################################

\subsubsection{Inviscid\slash Slip Wall}\label{SEC:BC:Wall:Slip}

The inviscid wall boundary condition, sometimes referred to as a slip wall, is
used to model a frictionless solid-wall surface by enforcing the flow velocity
and primitive variable gradients to be tangent to the boundary surface.
%
The details provided here also apply to the symmetry boundary condition
(\ref{SEC:BC:Numerical:Symmetry}) because its formulation happens to be
identical to a slip wall boundary condition.

%%%%%%%%%%%%%%%%%%%%%%%%%%%%%%%%%%%%%%%%%%%%%%%%%%%%%%%%%%%%%%%%%%%%%%%%%%%%%%%%

For this boundary condition, there is no momentum or heat flux through the
boundary so the thermodynamic properties at the wall are copied from the
interior:
%
\begin{alignat*}{1}
  \rho_\subR & = \rho_\subL \\
  T_\subR & = T_\subL
\end{alignat*}
%
To enforce tangency of the velocity at the wall, the wall velocity,
$\vec{u}_\subR$, is computed by taking the interior velocity vector,
$\vec{u}_\subL$, and subtracting out the normal components of the velocity
vector, $(\vec{u}_\subL\cdot\normal)\normal$, resulting in the tangential
velocity at the surface:
%
\begin{equation}
  \vec{u}_\subR = \vec{u}_\subL - \left(\vec{u}_\subL\cdot\normal\right)\normal
  \label{EQ:BC:TangentialVelocity}
\end{equation}
%
For those interested, \eref{EQ:BC:TangentialVelocity} can be
\quoterm{simplified} into a triple vector product as:
%
\begin{equation*}
  \vec{u}_\subR = \normal\times\left(\vec{u}_\subL\times\normal\right)
\end{equation*}
 
%The cross product $\vec{u}\cross\normal$ first creates a vector that is
%perpendicular to the plane created by the velocity and normal vectors.
%%
%Taking the cross product of the normal with the result of the first cross
%product then produces a vector that is perpendicular to the normal and lies in
%the plane created by the velocity and normal vectors, which is just the vector
%tangent to the surface in the direction of the original velocity vector.

%%%%%%%%%%%%%%%%%%%%%%%%%%%%%%%%%%%%%%%%%%%%%%%%%%%%%%%%%%%%%%%%%%%%%%%%%%%%%%%%

For this boundary condition, the wall-normal component of the primitive variable
gradients should be zero at the wall.
%
To accomplish this, we repurpose \eref{EQ:BC:TangentialVelocity} to compute
tangential gradient vectors for each primitive variable at the surface given the
interior gradient vector using
%
\begin{equation*}
  \vec{\varphi}_\subR = \vec{\varphi}_\subL
                   - \left(\vec{\varphi}_\subL\cdot\normal\right)\normal
\end{equation*}
%
where $\vec{\varphi}$ is a stand-in for the gradient of one of the primitive
variables, \thatis \[\vec{\varphi} = \nabla\rho, \nabla u, \nabla v, \nabla w,
\text{ or } \nabla T\]

%###############################################################################



%###############################################################################

\subsubsection{Adiabatic Viscous\slash No-Slip Wall}
\label{SEC:BC:Wall:Adiabatic}

The adiabatic viscous\slash no-slip wall boundary condition enforces a
prescribed wall velocity at the boundary surface with zero heat flux through the
boundary, \thatis $\partial T/\partial \normal = 0$.
%
Note, this boundary condition defaults to a stationary wall with zero wall
velocity, but this can be overridden by specifying non-zero values for any of
the boundary condition velocity components.
%
The thermodynamic properties for this boundary condition are just copied from
the interior:
%
\begin{alignat*}{1}
  \rho_\subR & = \rho_\subL \\
  T_\subR & = T_\subL
\end{alignat*}

%%%%%%%%%%%%%%%%%%%%%%%%%%%%%%%%%%%%%%%%%%%%%%%%%%%%%%%%%%%%%%%%%%%%%%%%%%%%%%%%

For most use cases, the prescribed wall velocity is simply zero, \forexample a
stationary flat plate.
%
However, there are certain situations where prescribing a non-zero wall velocity
is needed; a canonical example would be the classic lid-driven cavity problem,
and a more real-world example would be a spinning hub in a turbomachinery
problem.
%
If the prescribed wall velocity is not rotational, the following components of
\BCinptext{} directly set the wall velocity for the boundary condition:
%
\begin{InputExample}
  \item[{\BCinpBC[nmlBC::BCInput/VelocityX]{VelocityX}}]
  \item[{\BCinpBC[nmlBC::BCInput/VelocityY]{VelocityY}}]
  \item[{\BCinpBC[nmlBC::BCInput/VelocityZ]{VelocityZ}}]
\end{InputExample}
%
If the prescribed wall velocity is rotational, the input array
%
\begin{InputExample}
\item[{\BCinpBC[nmlBC::BCInput/RotationRate]{Rotation\_Rate(1:3)}[Rotation\_Rate]}]
\end{InputExample}
%
is used to specify the rotation rate, $\vec{\Omega}$, (in RPMs) for all wall
points in the boundary condition.
%
%The prescribed wall velocity at each wall point depends on the radius
%of the wall point to the axis of rotation, $\vec{r}$, and is given by
%(assuming implicit unit conversions):
%
The prescribed wall velocity at each wall point depends on the radius,
$\vec{r}$, from the axis of rotation to the wall point.
%
Since the axis of rotation for rotating boundaries originates from the
specified 
\inpBC[nmlBC::BCInput/Periodic/RotationCenter]{RotationCenter}%&
      [Periodic\%RotationCenter][],
$\vec{x}_c$, we can also write this in terms of the position vector.
%
The relation is given by (assuming implicit unit conversions):
%
\begin{equation*}
  \vec{u}_\subR = \vec{\Omega} \times \vec{r}
                = \vec{\Omega} \times \left( \vec{x} - \vec{x}_c \right)
\end{equation*}
%
Note, this is different from prescribing a non-rotational wall velocity where
all wall points are set to the same velocity.

%%%%%%%%%%%%%%%%%%%%%%%%%%%%%%%%%%%%%%%%%%%%%%%%%%%%%%%%%%%%%%%%%%%%%%%%%%%%%%%%

If using a moving reference frame, the input
%
\begin{InputExample}
  \item[{\BCinpBC[nmlBC::BCInput/IsSpecifiedInMRF]{Is\_Specified\_In\_MRF}}]
\end{InputExample}
%
can be used to simplify the boundary condition inputs for cases involving
multiple solid wall boundaries with velocities dependent on which reference
frame the boundary is viewed from.
%
An example of this would be a turbomachinery fan containing a spinning hub and
rotors --- with a rotation rate of $\vec{\Omega}$ --- inside a stationary case,
which is how this problem is viewed from the inertial\slash fixed\slash absolute
reference frame.
%
However, when viewed from the rotating reference frame, the hub and rotor are
stationary and the case is now rotating at $(-\vec{\Omega})$.
%
Note, the prescribed wall velocity components are specified as the absolute
velocity components.

%%%%%%%%%%%%%%%%%%%%%%%%%%%%%%%%%%%%%%%%%%%%%%%%%%%%%%%%%%%%%%%%%%%%%%%%%%%%%%%%

To illustrate the use of \BCinptext{Is\_Specified\_In\_MRF}, let us continue
with this turbomachinery example of a spinning hub\slash rotor surrounded by a
stationary case, and our simulation is using a rotating reference frame with
the same rotation rate as the spinning hub\slash rotor.
%
Let \BCinptext[1]{} be the boundary condition inputs for the casing,
and \BCinptext[2]{} be for the hub\slash rotor.
%
Using the default values of zero for the three velocity components, the boundary
conditions for this case can be provided as:
%
\begin{InputExample}
  \item[{\BCinptext[1]{Is\_Specified\_In\_MRF}} = \fls]
  \item[{\BCinptext[2]{Is\_Specified\_In\_MRF}} = \tru]
\end{InputExample}
%
One benefit from this approach is that the boundary condition wall velocities
will automatically adjust if the rotation rate of the rotating reference frame
is changed.

%%%%%%%%%%%%%%%%%%%%%%%%%%%%%%%%%%%%%%%%%%%%%%%%%%%%%%%%%%%%%%%%%%%%%%%%%%%%%%%%

For this boundary condition, the temperature gradient is the only right gradient
state treated differently from the other flow-through boundary conditions
(\eref{EQ:BC:FlowThroughGradient}).
%
The adiabatic condition prescribing zero heat flux through the boundary requires
the temperature gradient to be tangential to the boundary surface.
%
This is exactly the same as the temperature gradient boundary condition for the
slip wall in \sref{SEC:BC:Wall:Slip}, which is
%
\begin{equation*}
  \nabla T_\subR = \nabla T_\subL
                 - \left(\nabla T_\subL\cdot\normal\right)\normal
\end{equation*}

%%%%%%%%%%%%%%%%%%%%%%%%%%%%%%%%%%%%%%%%%%%%%%%%%%%%%%%%%%%%%%%%%%%%%%%%%%%%%%%%

%\textcolor{red}{
%%
%  The viscous wall condition, sometimes referred to as a no-slip wall
%  condition, enforces the prescribed wall tangential velocity along the
%  boundary.
%%
%  The wall normal velocity is set to zero, resulting in no flux through the
%  boundary.
%%
%  For a non-moving viscous wall, all components of the wall velocity are zero. 
%%
%  The adiabatic condition also enforces no heat flux through the boundary, so
%  there will be no wall-normal temperature gradient.
%%
%}

%###############################################################################



%###############################################################################

\subsubsection{Isothermal Viscous\slash No-Slip Wall}
\label{SEC:BC:Wall:Isothermal}

The isothermal viscous\slash no-slip wall boundary condition enforces a
prescribed wall velocity and temperature at the wall boundary surface.
%
The isothermal no-slip wall boundary condition is very similar to the adiabatic
no-slip wall in \sref{SEC:BC:Wall:Adiabatic}, with the only difference between
them being how the wall temperature and temperature gradients are set.
%
Therefore, the remainder of this section only provides details for what is done
differently from the adiabatic wall boundary condition.
%
Particularly, all of the details in \sref{SEC:BC:Wall:Adiabatic} on how the
prescribed wall velocity is determined are exactly the same for isothermal
walls.

%%%%%%%%%%%%%%%%%%%%%%%%%%%%%%%%%%%%%%%%%%%%%%%%%%%%%%%%%%%%%%%%%%%%%%%%%%%%%%%%

The backbone of the isothermal wall boundary condition is prescribing a fixed
wall temperature for a solid boundary surface.
%
The prescribed wall temperature for the boundary condition, $T_\subBC$, is
specified using the input
%\BCinptext{Wall\_Temp}.
\BCinpBC[nmlBC::BCInput/WallTemp]{Wall\_Temp}.
%
For the rest of the thermodynamic gas properties, pressure is taken from the
interior, and then the ideal gas law is used to compute the density at the wall,
resulting in the following right state variables:
%
\begin{alignat*}{1}
  T_\subR & = T_\subBC \\
  p_\subR & = p_\subL \\
  \rho_\subR & = \frac{p_\subR}{\rgas T_\subR}
\end{alignat*}

%%%%%%%%%%%%%%%%%%%%%%%%%%%%%%%%%%%%%%%%%%%%%%%%%%%%%%%%%%%%%%%%%%%%%%%%%%%%%%%%

Finally, unlike for adiabatic walls, the isothermal wall temperature gradient is
just copied from the interior,
%
\begin{equation*}
  \nabla T_\subR = \nabla T_\subL
\end{equation*}
%
Since this is also done for the gradients of all of the other variables, the
right gradient states for the isothermal wall are just the standard
\quoterm{flow-through} gradient boundary condition from
\eref{EQ:BC:FlowThroughGradient} used by most boundary conditions.

\devNotes{If default flow conditions are found for an isothermal wall, it will
use the \inptext{ovar\_mod} variable \inptext{iso\_wall\_temp} to set the wall
temperature to 300 K. This probably needs to be changed to an error+abort. This
dev-note is also in \sref{SEC:BC:Specify:FullState}.}

%###############################################################################



%###############################################################################
%###############################################################################
%###############################################################################
%###############################################################################
%###############################################################################



%###############################################################################

\subsection{Numerical Boundaries}\label{SEC:BC:Numerical}

%###############################################################################

\subsubsection{Symmetry Boundary Condition}
\label{SEC:BC:Numerical:Symmetry}

This boundary condition is generally used to reduce the size of the overall
computational domain by taking advantage of flow symmetry for a given problem.
%
For example, instead of running a simulation of a full airplane, a symmetry
boundary can be created along the centerline of the airplane geometry so the
simulation is only performed on half of the airplane.
%
Although intended for different uses, this boundary condition is equivalent to
the inviscid\slash slip wall boundary condition from \sref{SEC:BC:Wall:Slip}.

%%%%%%%%%%%%%%%%%%%%%%%%%%%%%%%%%%%%%%%%%%%%%%%%%%%%%%%%%%%%%%%%%%%%%%%%%%%%%%%%

\NOTE It should be warned that using symmetry boundary conditions to take
advantage of flow symmetry is generally not recommended for quality LES.
%
The mean flow may be symmetric over long time-averages, but the instantaneous
turbulent flow structures are \emph{not} symmetric.

%###############################################################################



%###############################################################################

\subsubsection{Periodic Coupled Boundaries}
\label{SEC:BC:Numerical:Periodic}

%###############################################################################

\devNotes{%
%
  \\\indent
%
  Currently, the translational vector is specified in grid scaled units.
  %
  This needs to be changed to grid units.
%
  }

%###############################################################################

Periodic coupled boundaries permit the flow exiting one boundary of the
computational domain to enter through a different boundary, and vice versa.
%
These two boundaries form a coupled pair of periodic boundary conditions.
%
Periodic boundaries fall into two categories of periodicity: translational and
rotational. A pair of periodic boundaries can be either translational,
rotational, or a combination of both.
%
Translational periodicity is quite useful for problems where the flow in one or
more directions can be treated as infinite, \forexample channel flow; details
and examples are discussed in
\sref{SEC:BC:Numerical:Periodic:Inputs:Translated}.
%
Rotational periodicity is useful for problems containing rotational symmetries,
\forexample a single blade from a turbomachinery rotor; details and examples are
discussed in \sref{SEC:BC:Numerical:Periodic:Inputs:Rotated}.

%%%%%%%%%%%%%%%%%%%%%%%%%%%%%%%%%%%%%%%%%%%%%%%%%%%%%%%%%%%%%%%%%%%%%%%%%%%%%%%%

We first start with some preliminary definitions and notation that will be used
for periodic coupled boundary conditions.
%
For all examples in this section, we will be working with a pair of periodic
boundary conditions that are specified in the input file using \BCinptext[i]{}
and \BCinptext[j]{}.
%
For the remainder of this section, every instance of \BCinptext[i]{} will have
the following two specified values
%
\begin{InputExample}
  \item[{\BCinptext[i]{Name = "Face Side A"}}]
  \item[{\BCinptext[i]{Set\_BC\_Value = 81}}]
\end{InputExample}
%
and similarly, \BCinptext[j] will have the following two specified values
%
\begin{InputExample}
  \item[{\BCinptext[j]{Name = "Periodic-2"}}]
  \item[{\BCinptext[j]{Set\_BC\_Value = 81}}]
\end{InputExample}
%
Note, the strings used here for the \inptext{Name} components of each boundary
condition have just been randomly created to illustrate that there is no
restriction that they be coordinated in any way.
%
In practice, it is recommended that they are coordinated just to make things
easier to understand, say something like \inptext{"Periodic-1A"} and
\inptext{"Periodic-1B"} if there might be more than one pair of periodic
boundary conditions.

%###############################################################################

\subsubsubsection{Things to Consider when Using Periodic Boundary Conditions}

\begin{itemize}
  %
  \item
    After all translation and rotation operations are performed on one side of a
    pair of coupled periodic boundaries, the linear grid components (linear cell
    vertices and faces) of the two boundaries need to result in a 1-to-1 match.
    %
    The algorithm for finding all of the matching pairs of vertices between two
    periodic boundaries works by computing the distance from the current vertex
    from the left boundary to all the vertices on the right boundary that have
    not yet been matched to a left vertex.
    %
    The matching right vertex is simply determined to be the one with the
    minimum distance to the current left vertex.
    %
    If the minimum distance is greater than 0.001
    \key[key:GridScaledUnits]{grid scaled units}[grid scaled units]
    (\thatis meters), an error will occur.

    For a pair of matching linear cell faces on periodic boundaries, the
    geometries for both faces of the pair need to match (\thatis both
    triangles, or both quadrilaterals), and the set of grid vertices for both
    faces need to create a matching set.
    %
    For example, let us say we have face $a$ defined by vertices $[1,2,3]$, and
    it was found to be a periodic match for face $b$ defined by vertices
    $[143,22,200]$.
    %
    For this to be a valid match, vertex 1 needs to match to one of the three
    vertices for face $b$; for this example we will say it matches to vertex
    200.
    %
    Vertex 2 then needs to be a match to either vertex 143 or 22; we will say 22
    for this example.
    %
    That means vertex 3 \emph{must} be a match to vertex 143, or else this is
    not a valid pair of matching faces.
    %
    Finally, the order of the matched vertices is not important here since after
    the faces are confirmed to be a valid match, the ordering will be used to
    find the relative rotation between the two faces.
  %
  \item
    Provided the previous bullet is satisfied such that the linear grid
    components of a pair of periodic boundaries are found to be an exact
    1-to-1 match, there are no other requirements on the topology, shape,
    or orientation of the boundaries.
    %
    Periodic boundaries can be high-order, curved surfaces, and they are not
    required to be planar or aligned with any coordinate direction\slash axis.
    %
    For example, a two-dimensional infinite turbine blade cascade could be
    simulated using only a single blade with translationally periodic boundaries
    along the curved geometric mid-passage lines above\slash below the blade.
  %
  \item
    For LES, make sure the distance between a pair of coupled
    periodic boundaries are separated by a distance that is sufficient enough to
    allow for turbulent structures to breakdown\slash dissipate\slash cascade
    into finer turbulent length scales along the direction of the periodicity.
    %
    If a pair of coupled periodic boundaries are too close to each other, the
    turbulent length scale will get artificially constrained in the periodic
    direction.
    %
    This will cause energy to get trapped in the larger\slash macro vortical
    structures, ultimately preventing a truly three-dimensional turbulent flow.
    
   %\FIXME{\\
   %This discusses breakdown, but don't periodic boundaries inhibit
   %\emph{growth} of turbulent structures?\\
   %Recommendation for spanwise width for LES?
   %Perhaps some multiple of the expected turbulent length scale?
   %}
  %
  \item
    Similar to the previous bullet, be careful with over-prescribing
    a simulation with too many periodic boundaries in such a way that the
    periodic boundary conditions change the flow dynamics.
    %
    The use of periodic boundary conditions is actually creating an infinite
    array of identical coupled simulations that are interacting with each other.
    %
    If the effect of this coupling is not well thought out, it can have
    unforeseen consequences ranging from unstable simulations that can be hard
    to diagnose the source of the instabilities, to questions on the validity
    \andor accuracy of a simulation and any conclusions that may have resulted.
     
    As mentioned in the previous bullet, having periodic boundaries too close to
    each other is an example of this, but another example is the fully periodic
    isentropic vortex problem\cite{Spiegel_2015}.
    %
    For this problem, 
    1) the Euler equations provide no physical dissipation;
    2) the high-order FR scheme provides minimal numerical dissipation; and
    3) the fully periodic boundaries prevent both numerical and machine
       roundoff errors from escaping the system.
    %
    With no means to dissipate the error or let it leave the system, all of the
    numerical and roundoff errors created by the simulation become trapped, and
    eventually these errors grow large enough to negatively effect the
    simulations with unstable\slash diverging solutions.
    %
    Replacing the periodic boundaries with characteristic boundary conditions
    completely eliminates all instabilities allowing us to demonstrate that
    the FR method achieves super accuracy for this problem.
    %
    Further details about this can be found in \linecite{Spiegel_2015}.
  %
\end{itemize}

%###############################################################################

\subsubsubsection{Specifying Coupled Periodic Boundary Conditions}
\label{SEC:BC:Numerical:Periodic:Inputs}

%###############################################################################

When specifying periodic boundary conditions, some input parameters are
commonly used with all such boundaries, and others are specific to
translational or rotational periodicity.
%
We discuss each in the sections that follow.

%###############################################################################

\FourSubsection{Inputs Common to All Periodic Boundaries}
\label{SEC:BC:Numerical:Periodic:Inputs:Common}

%This section discusses the inputs for specifying a pair of coupled periodic
%boundaries that are common to all periodic boundaries whether the periodicity
%is translational or rotational.
%
This section discusses inputs that can be used with both translational or
rotational periodic boundaries.
%
The \inptext{Periodic} subcomponent of \BCinpBC[nmlBC::BCInput]{} is used
to provide the information that will allow GFR to group two boundary conditions
into a matching pair of periodic boundaries, which in this example will be
\BCinptext[i]{} and \BCinptext[j]{}.
%
Both sides of a periodic pair must provide their respective information in the
input file.
%
%As you will see, the information required by opposing sides is the reverse\slash
%negative of each other, and could therefore probably be streamlined to reduce
%redundancies.
%%
%However, there was intentionality in how this was set up to help aid in error
%checking and making sure the user understands what they are doing in setting up
%periodic boundary conditions.
%
As you will see, the information required by opposing sides is the reverse\slash
negative of each other, and is used to aid error checking of the periodic
boundary setup.

%%%%%%%%%%%%%%%%%%%%%%%%%%%%%%%%%%%%%%%%%%%%%%%%%%%%%%%%%%%%%%%%%%%%%%%%%%%%%%%%

The two inputs used to define a matching periodic pair are
%
\begin{InputExample}
  \item[{\BCinptext[i]{Periodic\%Matching\_BC\_Name}}]
  \item[{\BCinptext[i]{Periodic\%Matching\_BC\_Idx}}]
\end{InputExample}
%
The values for these two inputs need to provide the respective information from
the opposing matching periodic boundary condition.
%
Using boundary conditions $i$ and $j$ defined at the beginning of
\sref{SEC:BC:Numerical:Periodic}, the inputs for boundary condition $i$ would
look like
%
\begin{InputExample}
  \item[{\BCinptext[i]{Name = "Face Side A"}}]
  \item[{\BCinptext[i]{Set\_BC\_Value = 81}}]
  \item[{\BCinptext[i]{Periodic\%Matching\_BC\_Name = "Periodic-2"}}]
  \item[{\BCinptext[i]{Periodic\%Matching\_BC\_Idx = }j}]
\end{InputExample}
%
and the inputs for boundary condition $j$ would look like
%
\begin{InputExample}
  \item[{\BCinptext[j]{Name = "Periodic-2"}}]
  \item[{\BCinptext[j]{Set\_BC\_Value = 81}}]
  \item[{\BCinptext[j]{Periodic\%Matching\_BC\_Name = "Face Side A"}}]
  \item[{\BCinptext[j]{Periodic\%Matching\_BC\_Idx = }i}]
\end{InputExample}
%
The remaining parts of \BCinptext{Periodic} are used to define the physical
orientation between the pair of periodic boundaries, and will be covered in the
next two sections.

%###############################################################################



%###############################################################################

\FourSubsection{Translated Periodic Boundary Condition}
\label{SEC:BC:Numerical:Periodic:Inputs:Translated}

Translated periodic boundary conditions are useful for problems where the flow
in one or more directions can be treated as infinite.
%
Examples primarily consist of two-dimensional geometries
(\forexample an airfoil)
that are converted to a three-dimensional body with an infinite span in the
third dimension.
%
Another example would be a 3D channel flow with translated periodic boundary
conditions in both the span-wise \emph{and} streamwise directions.
%
The periodicity in the streamwise direction allows for the flow to recirculate
indefinitely until it becomes a fully developed channel flow.
%
However, the streamwise periodicity requires a body force to drive the
streamwise bulk flow in the absence of a pressure differential generated by
traditional inflow\slash outflow boundary conditions.

%%%%%%%%%%%%%%%%%%%%%%%%%%%%%%%%%%%%%%%%%%%%%%%%%%%%%%%%%%%%%%%%%%%%%%%%%%%%%%%%

For each periodic boundary condition, the input vector
%
\begin{InputExample}
  \item[{\BCinpBC[nmlBC::BCInput/Periodic/Translation][i]%&
                 {Periodic\%Translation(1:3)}%
                 [Periodic\%Translation]}]
\end{InputExample}
%
provides the translation vector (in grid
units) that transforms the current boundary surface to align it exactly with its
corresponding periodic pair.
%
For a given periodic pair with translational periodicity, the following relation
must be satisfied:
%
\begin{InputExample}
  \item[{\BCinptext[i]{Periodic\%Translation(:) = -}
         \BCinptext[j]{Periodic\%Translation(:)}}]
\end{InputExample}

%%%%%%%%%%%%%%%%%%%%%%%%%%%%%%%%%%%%%%%%%%%%%%%%%%%%%%%%%%%%%%%%%%%%%%%%%%%%%%%%

For example, say you have an S-duct that is 5.0 inches long and drops 1.0 inch
vertically, and your grid file for this geometry is also provided in inches.
%
If you wanted to make this problem streamwise periodic, your inputs for boundary
condition $i$ would look like
%
\begin{InputExample}
  \item[{\BCinptext[i]{Periodic\%Translation(1:3) = [ 5.0, 1.0, 0.0]}}]
 %\item[{\BCinptext[i]{Periodic\%Translation(1) = 5.0}}]
 %\item[{\BCinptext[i]{Periodic\%Translation(2) = 1.0}}]
 %\item[{\BCinptext[i]{Periodic\%Translation(3) = 0.0}}]
\end{InputExample}
%
and the inputs for boundary condition $j$ would look like
%
\begin{InputExample}
  \item[{\BCinptext[j]{Periodic\%Translation(1:3) = [-5.0, -1.0, 0.0]}}]
 %\item[{\BCinptext[j]{Periodic\%Translation(1) = -5.0}}]
 %\item[{\BCinptext[j]{Periodic\%Translation(2) = -1.0}}]
 %\item[{\BCinptext[j]{Periodic\%Translation(3) = 0.0}}]
\end{InputExample}
%
Of course, an S-duct is generally not a problem that you would want to make
streamwise periodic.
%
It is merely used here to illustrate an example of the
periodic translational vector having components in multiple directions.

%###############################################################################

The elements on the boundary faces must be identical in order to have one-to-one
connectivity between cell faces.
%
When using translated periodic boundary conditions for large-eddy simulations,
one must make sure the extent of the periodic direction is long enough as to not
artificially constrain the turbulent length scale in the periodic direction.

%%%%%%%%%%%%%%%%%%%%%%%%%%%%%%%%%%%%%%%%%%%%%%%%%%%%%%%%%%%%%%%%%%%%%%%%%%%%%%%%



%###############################################################################

\FourSubsection{Rotated Periodic Boundary Condition}
\label{SEC:BC:Numerical:Periodic:Inputs:Rotated}

Rotated periodic boundary conditions are useful for problems containing
rotational symmetry.
%
This is particularly prevalent in turbomachinery flows where simulations of a
single periodic rotor blade can save significant computational expense compared
to simulations of the full 360-degree rotor wheel with dozens of rotor blades.
%
There are three subcomponent arrays of \BCinptext[i]{Periodic} that can be used
to specify the rotational transformation between a pair of rotated periodic
boundary conditions:
%
\begin{InputExample}
  \item[{\BCinpBC[nmlBC::BCInput/Periodic/RotationCenter][i]%&
                 {Periodic\%RotationCenter(1:3)}%
                 [Periodic\%RotationCenter]}]
  \item[{\BCinpBC[nmlBC::BCInput/Periodic/RotationAngle][i]%&
                 {Periodic\%RotationAngle(1:3)}%
                 [Periodic\%RotationAngle]}]
  \item[{\BCinpBC[nmlBC::BCInput/Periodic/RotationAxis][i]%&
                 {Periodic\%RotationAxis(1:4)}%
                 [Periodic\%RotationAxis]}]
\end{InputExample}
%
%One may use either \BCinptext[i]{Periodic\%RotationAngle} or
%\BCinptext[i]{Periodic\%RotationAxis} to define rotated periodic boundary
%conditions, but not both.
%
One may use either \inptext{RotationAngle} or \inptext{RotationAxis} to define
rotated periodic boundary conditions, but not both.
%
If both are specified, \inptext{RotationAxis} will be ignored and
\inptext{RotationAngle} will be used.
%
%For both methods, the input vector \BCinptext[i]{Periodic\%RotationCenter(1:3)}
%gives the rotation origin (in \key[key:GridUnits]{grid units}) used
%for defining the rotation angles between rotated periodic boundary conditions.
%
For both methods, the input vector \inptext{RotationCenter(1:3)} gives the
coordinates for the rotation origin (in \key[key:GridUnits]{grid units}) used
for defining the rotation angles between rotated periodic boundary conditions.

%%%%%%%%%%%%%%%%%%%%%%%%%%%%%%%%%%%%%%%%%%%%%%%%%%%%%%%%%%%%%%%%%%%%%%%%%%%%%%%%

The first way to specify the rotational transformation between a pair of rotated
periodic boundary conditions is through the input vector
\inptext{RotationAngle(1:3)}.
%\BCinptext[i]{Periodic\%RotationAngle(1:3)}.
%
This vector gives the rotation angles (in degrees) about each of the coordinate
axes from the current periodic boundary interface to the connecting periodic
boundary interface.
%
If rotating about more than one axis, the rotation is performed first about the
$x$-axis, then the $y$-axis, and finally the $z$-axis.
%
For a given periodic pair with rotational periodicity, the following relation
needs to be true:
%
\begin{InputExample}
  \item[\small{\BCinptext[i]{Periodic\%RotationAngle(:) = -}
               \BCinptext[j]{Periodic\%RotationAngle(:)}}]
\end{InputExample}

%%%%%%%%%%%%%%%%%%%%%%%%%%%%%%%%%%%%%%%%%%%%%%%%%%%%%%%%%%%%%%%%%%%%%%%%%%%%%%%%

The second way to specify the rotational transformation between a pair of
rotated periodic boundary conditions is through the input vector
\inptext{RotationAxis(1:4)}.
%\BCinptext[i]{Periodic\%RotationAxis(1:4)}.
%
The first 3 values of this array define a vector representing the axis of
rotation for this rotated periodic boundary condition, and the last value of the
array defines the rotation angle (in degrees) about this axis.
%
The magnitude of the vector defining the axis of rotation is not used; the sole
purpose of this vector is to define the direction of the rotation axis.
%
For a given periodic pair with rotational periodicity, one of the two following
sets of relations needs to be true:
%
\begin{InputExample}
  \item[\small{\BCinptext[i]{Periodic\%RotationAxis(1:3) = }
               \BCinptext[j]{Periodic\%RotationAxis(1:3)}}]
  \item[\small{\BCinptext[i]{Periodic\%RotationAxis(4) = -}
               \BCinptext[j]{Periodic\%RotationAxis(4)}}]
\end{InputExample}
%
or
%
\begin{InputExample}
  \item[\small{\BCinptext[i]{Periodic\%RotationAxis(1:3) = -}
               \BCinptext[j]{Periodic\%RotationAxis(1:3)}}]
  \item[\small{\BCinptext[i]{Periodic\%RotationAxis(4) = }
               \BCinptext[j]{Periodic\%RotationAxis(4)}}]
\end{InputExample}

%%%%%%%%%%%%%%%%%%%%%%%%%%%%%%%%%%%%%%%%%%%%%%%%%%%%%%%%%%%%%%%%%%%%%%%%%%%%%%%%

For example, say you have a rotor fan with 15 blades with the rotation axis
aligned with the Cartesian $x$-axis, and you would like to run a simulation of
just a single blade using rotational periodicity.
%
Using \inptext{RotationAngle}, your inputs for boundary condition $i$ would
look like
%
\begin{InputExample}
  \item[{\BCinptext[i]{Periodic\%RotationAngle(1:3) = [ 24.0, 0.0, 0.0]}}]
\end{InputExample}
%
and the inputs for boundary condition $j$ would look like
%
\begin{InputExample}
  \item[{\BCinptext[j]{Periodic\%RotationAngle(1:3) = [-24.0, 0.0, 0.0]}}]
\end{InputExample}
%
If you wanted to use \inptext{RotationAxis} instead, your inputs for boundary
condition $i$ would look like
%
\begin{InputExample}
  \item[{\BCinptext[i]{Periodic\%RotationAxis(1:4) = [ 1.0, 0.0, 0.0, 24.0]}}]
\end{InputExample}
%
and the inputs for boundary condition $j$ would look like
%
\begin{InputExample}
  \item[{\BCinptext[j]{Periodic\%RotationAxis(1:4) = [ 1.0, 0.0, 0.0,-24.0]}}]
\end{InputExample}



%###############################################################################

\subsubsubsection{Past Examples of Periodic Boundary Conditions}
\label{SEC:BC:Numerical:Periodic:Examples}

These are examples taken from GFR input files that have been used in the past to
run simulations for problems including periodic boundaries.
%
Note that the array indices used for \inptext{BC\_Input} are simply copied over
from the input files, meaning each index is relative to the other boundary
conditions that also existed in the input file it was taken from.
%
Part of the intent behind keeping these array indices as they originally were in
their respective input file is to demonstrate that the previously used variable
indices of $i$ and $j$ are can represent any array element of
\inptext{BC\_Input} provided that they link to each other when they form a pair
of periodic boundaries.
%
\begin{itemize}
  \item Streamwise and spanwise periodic channel flow
   %\begin{nmldescitem}[escapechar=@]
   %\begin{nmldescitem}[fontupper=\ttfamily\footnotesize,
    \begin{nmldescitem}[fontupper=\ttfamily\scriptsize,
                        listing options={escapeinside={!@}{!@}}]
&nml_bc
  BC_Input(1)%Name = "INFLOW"
  BC_Input(1)%Set_BC_Value = 81
  BC_Input(1)%Periodic%Matching_BC_Name = "OUTFLOW"
  BC_Input(1)%Periodic%Matching_BC_Idx = 2
  BC_Input(1)%Periodic%Translation(1) = 0.036576 !@\vspace{5pt}!@
  BC_Input(2)%Name = "OUTFLOW"
  BC_Input(2)%Set_BC_Value = 81
  BC_Input(2)%Periodic%Matching_BC_Name = "INFLOW"
  BC_Input(2)%Periodic%Matching_BC_Idx = 1
  BC_Input(2)%Periodic%Translation(1) = -0.036576 !@\vspace{5pt}!@
  BC_Input(5)%Name = "ZMIN"
  BC_Input(5)%Set_BC_Value = 81
  BC_Input(5)%Periodic%Matching_BC_Name = "ZMAX"
  BC_Input(5)%Periodic%Matching_BC_Idx = 6
  BC_Input(5)%Periodic%Translation(3) = 0.018288 !@\vspace{5pt}!@
  BC_Input(6)%Name = "ZMAX"
  BC_Input(6)%Set_BC_Value = 81
  BC_Input(6)%Periodic%Matching_BC_Name = "ZMIN"
  BC_Input(6)%Periodic%Matching_BC_Idx = 5
  BC_Input(6)%Periodic%Translation(3) = -0.018288
/
    \end{nmldescitem}
   %\FINALFIX{\newpage}
  \item Single rotor blade from spacecraft cabin ventilation
    fan\cite{Spiegel_2025}
    \begin{itemize}[before=\small]
      \item Rotor consists of nine blades spinning about the $x$-axis resulting
        in a rotationally periodic domain that is rotated 40\textdegree{} about
        the $x$-axis.
      \item This grid has two pairs of rotationally periodic boundaries: one for
        the primary computational domain, and a second for a sponge region
        appended onto the regular outflow boundary to damp out non-physical
        pressure waves from reflecting off the outflow boundary.
        The periodic boundary is split in this way for convenience, but could
        be combined.
    \end{itemize}
   %\begin{nmldescitem}[escapechar=@]
   %\begin{nmldescitem}[fontupper=\ttfamily\footnotesize,
    \begin{nmldescitem}[fontupper=\ttfamily\scriptsize,
                        listing options={escapeinside={!@}{!@}}]
&nml_bc
  BC_Input(8)%Name = "ThetaMin"
  BC_Input(8)%Set_BC_Value = 81
  BC_Input(8)%Periodic%Matching_BC_Name = "ThetaMax"
  BC_Input(8)%Periodic%Matching_BC_Idx = 9
  BC_Input(8)%Periodic%RotationAngle(1) = 40.0 !@\vspace{5pt}!@
  BC_Input(9)%Name = "ThetaMax"
  BC_Input(9)%Set_BC_Value = 81
  BC_Input(9)%Periodic%Matching_BC_Name = "ThetaMin"
  BC_Input(9)%Periodic%Matching_BC_Idx = 8
  BC_Input(9)%Periodic%RotationAngle(1) = -40.0 !@\vspace{5pt}!@
  BC_Input(10)%Name = "ThetaMinSponge"
  BC_Input(10)%Set_BC_Value = 81
  BC_Input(10)%Periodic%Matching_BC_Name = "ThetaMaxSponge"
  BC_Input(10)%Periodic%Matching_BC_Idx = 11
  BC_Input(10)%Periodic%RotationAngle(1) = 40.0 !@\vspace{5pt}!@
  BC_Input(11)%Name = "ThetaMaxSponge"
  BC_Input(11)%Set_BC_Value = 81
  BC_Input(11)%Periodic%Matching_BC_Name = "ThetaMinSponge"
  BC_Input(11)%Periodic%Matching_BC_Idx = 10
  BC_Input(11)%Periodic%RotationAngle(1) = -40.0
/
    \end{nmldescitem}
  \item Ceramic Matrix Composite 9\% thickness (CMC9) blade
    cascade\cite{Giel_2020} with spanwise and vertical
    periodicity\cite{Miki_2022}
   %\begin{nmldescitem}[escapechar=@]
   %\begin{nmldescitem}[fontupper=\ttfamily\footnotesize,
    \begin{nmldescitem}[fontupper=\ttfamily\scriptsize,
                        listing options={escapeinside={!@}{!@}}]
&nml_bc
  BC_Input(3)%Name = "periodic_yp13"
  BC_Input(3)%Set_BC_Value = 81
  BC_Input(3)%Periodic%Matching_BC_Name = "periodic_ym"
  BC_Input(3)%Periodic%Matching_BC_Idx = 4
  BC_Input(3)%Periodic%Translation(2) = -0.13 !@\vspace{5pt}!@
  BC_Input(4)%Name = "periodic_ym"
  BC_Input(4)%Set_BC_Value = 81
  BC_Input(4)%Periodic%Matching_BC_Name = "periodic_yp13"
  BC_Input(4)%Periodic%Matching_BC_Idx = 3
  BC_Input(4)%Periodic%Translation(2) = 0.13 !@\vspace{5pt}!@
  BC_Input(5)%Name = "periodic_zp013"
  BC_Input(5)%Set_BC_Value = 81
  BC_Input(5)%Periodic%Matching_BC_Name = "periodic_z0"
  BC_Input(5)%Periodic%Matching_BC_Idx = 6
  BC_Input(5)%Periodic%Translation(3) = -0.052 !@\vspace{5pt}!@
  BC_Input(6)%Name = "periodic_z0"
  BC_Input(6)%Set_BC_Value = 81
  BC_Input(6)%Periodic%Matching_BC_Name = "periodic_zp013"
  BC_Input(6)%Periodic%Matching_BC_Idx = 5
  BC_Input(6)%Periodic%Translation(3) = 0.052
/
    \end{nmldescitem}
  \item Taylor-Couette flow\cite{Wang_2021} with streamwise periodicity
   %\begin{nmldescitem}[escapechar=@]
   %\begin{nmldescitem}[fontupper=\ttfamily\footnotesize,
    \begin{nmldescitem}[fontupper=\ttfamily\scriptsize,
                        listing options={escapeinside={!@}{!@}}]
&nml_bc
  BC_Input(3)%Name = "Inflow"
  BC_Input(3)%Set_BC_Value = 81
  BC_Input(3)%Periodic%Matching_BC_Name = "Outflow"
  BC_Input(3)%Periodic%Matching_BC_Idx = 4
  BC_Input(3)%Periodic%Translation(1) = 6.283185307 !@\vspace{5pt}!@
  BC_Input(4)%Name = "Outflow"
  BC_Input(4)%Set_BC_Value = 81
  BC_Input(4)%Periodic%Matching_BC_Name = "Inflow"
  BC_Input(4)%Periodic%Matching_BC_Idx = 3
  BC_Input(4)%Periodic%Translation(1) = -6.283185307
/
    \end{nmldescitem}
\end{itemize}

%###############################################################################

%\begin{itemize}
%  \item Streamwise and spanwise periodic channel flow
%    \begin{InputExample}[before=\footnotesize]
%      %
%      \item[{\BCinptext[\texttt{1}]{Name = "INFLOW"}}]
%      \item[{\BCinptext[\texttt{1}]{Set\_BC\_Value = 81}}]
%      \item[{\BCinptext[\texttt{1}]{Periodic\%Matching\_BC\_Name = "OUTFLOW"}}]
%      \item[{\BCinptext[\texttt{1}]{Periodic\%Matching\_BC\_Idx = 2}}]
%      \item[{\BCinptext[\texttt{1}]{Periodic\%Translation(1) = 0.036576}}]
%      %
%      \item[{\BCinptext[\texttt{2}]{Name = "OUTFLOW"}}]\vspace{5pt}
%      \item[{\BCinptext[\texttt{2}]{Set\_BC\_Value = 81}}]
%      \item[{\BCinptext[\texttt{2}]{Periodic\%Matching\_BC\_Name = "INFLOW"}}]
%      \item[{\BCinptext[\texttt{2}]{Periodic\%Matching\_BC\_Idx = 1}}]
%      \item[{\BCinptext[\texttt{2}]{Periodic\%Translation(1) = -0.036576}}]
%      %
%      \item[{\BCinptext[\texttt{5}]{Name = "ZMIN"}}]\vspace{5pt}
%      \item[{\BCinptext[\texttt{5}]{Set\_BC\_Value = 81}}]
%      \item[{\BCinptext[\texttt{5}]{Periodic\%Matching\_BC\_Name = "ZMAX"}}]
%      \item[{\BCinptext[\texttt{5}]{Periodic\%Matching\_BC\_Idx = 6}}]
%      \item[{\BCinptext[\texttt{5}]{Periodic\%Translation(3) = 0.018288}}]
%      %
%      \item[{\BCinptext[\texttt{6}]{Name = "ZMAX"}}]\vspace{5pt}
%      \item[{\BCinptext[\texttt{6}]{Set\_BC\_Value = 81}}]
%      \item[{\BCinptext[\texttt{6}]{Periodic\%Matching\_BC\_Name = "ZMIN"}}]
%      \item[{\BCinptext[\texttt{6}]{Periodic\%Matching\_BC\_Idx = 5}}]
%      \item[{\BCinptext[\texttt{6}]{Periodic\%Translation(3) = -0.018288}}]
%      %
%    \end{InputExample}
%  \item Single rotor blade from spacecraft cabin ventilation
%    fan\cite{Spiegel_2025}
%    \begin{itemize}
%      \item Rotor consists of nine blades spinning about the $x$-axis resulting
%        in a rotationally periodic domain that is rotated 40\textdegree{} about
%        the $x$-axis.
%      \item This grid has two pairs of rotationally periodic boundaries: one for
%        the primary computational domain, and a second for a sponge region
%        appended onto the regular outflow boundary to damp out non-physical
%        pressure waves from reflecting off the outflow boundary.
%    \end{itemize}
%    \begin{InputExample}[before=\footnotesize]
%      %
%      \item[{\BCinptext[\texttt{8}]{Name = "ThetaMin"}}]
%      \item[{\BCinptext[\texttt{8}]{Set\_BC\_Value = 81}}]
%      \item[{\BCinptext[\texttt{8}]{Periodic\%Matching\_BC\_Name = "ThetaMax"}}]
%      \item[{\BCinptext[\texttt{8}]{Periodic\%Matching\_BC\_Idx = 9}}]
%      \item[{\BCinptext[\texttt{8}]{Periodic\%RotationAngle(1) = 40.0}}]
%      %
%      \item[{\BCinptext[\texttt{9}]{Name = "ThetaMax"}}]\vspace{5pt}
%      \item[{\BCinptext[\texttt{9}]{Set\_BC\_Value = 81}}]
%      \item[{\BCinptext[\texttt{9}]{Periodic\%Matching\_BC\_Name = "ThetaMin"}}]
%      \item[{\BCinptext[\texttt{9}]{Periodic\%Matching\_BC\_Idx = 8}}]
%      \item[{\BCinptext[\texttt{9}]{Periodic\%RotationAngle(1) = -40.0}}]
%      %
%      \item[{\BCinptext[\texttt{10}]{Name = "ThetaMinSponge"}}]\vspace{5pt}
%      \item[{\BCinptext[\texttt{10}]{Set\_BC\_Value = 81}}]
%      \item[{\BCinptext[\texttt{10}]{Periodic\%Matching\_BC\_Name = "ThetaMaxSponge"}}]
%      \item[{\BCinptext[\texttt{10}]{Periodic\%Matching\_BC\_Idx = 11}}]
%      \item[{\BCinptext[\texttt{10}]{Periodic\%RotationAngle(1) = 40.0}}]
%      %
%      \item[{\BCinptext[\texttt{11}]{Name = "ThetaMaxSponge"}}]\vspace{5pt}
%      \item[{\BCinptext[\texttt{11}]{Set\_BC\_Value = 81}}]
%      \item[{\BCinptext[\texttt{11}]{Periodic\%Matching\_BC\_Name = "ThetaMinSponge"}}]
%      \item[{\BCinptext[\texttt{11}]{Periodic\%Matching\_BC\_Idx = 10}}]
%      \item[{\BCinptext[\texttt{11}]{Periodic\%RotationAngle(1) = -40.0}}]
%      %
%    \end{InputExample}
%  \item Ceramic Matrix Composite 9\% thickness (CMC9) blade\cite{Giel_2020} with
%    spanwise and vertical periodicity\cite{Miki_2022}
%    \begin{InputExample}[before=\footnotesize]
%      %
%      \item[{\BCinptext[\texttt{3}]{Name = "periodic\_yp13"}}]
%      \item[{\BCinptext[\texttt{3}]{Set\_BC\_Value = 81}}]
%      \item[{\BCinptext[\texttt{3}]{Periodic\%Matching\_BC\_Name = "periodic\_ym"}}]
%      \item[{\BCinptext[\texttt{3}]{Periodic\%Matching\_BC\_Idx = 4}}]
%      \item[{\BCinptext[\texttt{3}]{Periodic\%Translation(2) = -0.13}}]
%      %
%      \item[{\BCinptext[\texttt{4}]{Name = "periodic\_ym"}}]\vspace{5pt}
%      \item[{\BCinptext[\texttt{4}]{Set\_BC\_Value = 81}}]
%      \item[{\BCinptext[\texttt{4}]{Periodic\%Matching\_BC\_Name = "periodic\_yp13"}}]
%      \item[{\BCinptext[\texttt{4}]{Periodic\%Matching\_BC\_Idx = 3}}]
%      \item[{\BCinptext[\texttt{4}]{Periodic\%Translation(2) = 0.13}}]
%      %
%      \item[{\BCinptext[\texttt{5}]{Name = "periodic\_zp013"}}]\vspace{5pt}
%      \item[{\BCinptext[\texttt{5}]{Set\_BC\_Value = 81}}]
%      \item[{\BCinptext[\texttt{5}]{Periodic\%Matching\_BC\_Name = "periodic\_z0"}}]
%      \item[{\BCinptext[\texttt{5}]{Periodic\%Matching\_BC\_Idx = 6}}]
%      \item[{\BCinptext[\texttt{5}]{Periodic\%Translation(3) = -0.052}}]
%      %
%      \item[{\BCinptext[\texttt{6}]{Name = "periodic\_z0"}}]\vspace{5pt}
%      \item[{\BCinptext[\texttt{6}]{Set\_BC\_Value = 81}}]
%      \item[{\BCinptext[\texttt{6}]{Periodic\%Matching\_BC\_Name = "periodic\_zp013"}}]
%      \item[{\BCinptext[\texttt{6}]{Periodic\%Matching\_BC\_Idx = 5}}]
%      \item[{\BCinptext[\texttt{6}]{Periodic\%Translation(3) = 0.052}}]
%      %
%    \end{InputExample}
%  \item Taylor-Couette flow\cite{Wang_2021} with streamwise periodicity
%    \begin{InputExample}[before=\footnotesize]
%      %
%      \item[{\BCinptext[\texttt{3}]{Name = "Inflow"}}]
%      \item[{\BCinptext[\texttt{3}]{Set\_BC\_Value = 81}}]
%      \item[{\BCinptext[\texttt{3}]{Periodic\%Matching\_BC\_Name = "Outflow"}}]
%      \item[{\BCinptext[\texttt{3}]{Periodic\%Matching\_BC\_Idx = 4}}]
%      \item[{\BCinptext[\texttt{3}]{Periodic\%Translation(1) = 6.283185307}}]
%      %
%      \item[{\BCinptext[\texttt{4}]{Name = "Outflow"}}]\vspace{5pt}
%      \item[{\BCinptext[\texttt{4}]{Set\_BC\_Value = 81}}]
%      \item[{\BCinptext[\texttt{4}]{Periodic\%Matching\_BC\_Name = "Inflow"}}]
%      \item[{\BCinptext[\texttt{4}]{Periodic\%Matching\_BC\_Idx = 3}}]
%      \item[{\BCinptext[\texttt{4}]{Periodic\%Translation(1) = -6.283185307}}]
%      %
%    \end{InputExample}
%\end{itemize}

%###############################################################################

\IDEA{Each example should include a schematic of the labeled boundary
surfaces to help visualize the periodicity.}

%###############################################################################



%###############################################################################
%###############################################################################
%###############################################################################
%###############################################################################
%###############################################################################



%###############################################################################
%###############################################################################
%###############################################################################
%###############################################################################
%###############################################################################



%###############################################################################

\endgroup

